% =====================================================================
%  Supplemental Material. Compile before main.tex (main reads supplement.aux).
% =====================================================================
\documentclass[aps,pra,onecolumn,notitlepage,superscriptaddress,nofootinbib,floatfix]{revtex4-2}
\usepackage{amsmath,amssymb,amsthm,bm}
\usepackage{graphicx}
\usepackage{booktabs}
\usepackage{longtable}
\usepackage{xcolor}
\usepackage{enumitem}
\usepackage[hidelinks,pdftitle={Supplemental Material: When does tail-risk optimization change a quantum gate?},
            pdfauthor={Sashreek Satapathy}]{hyperref}
\usepackage{xr-hyper}
\externaldocument[M-]{main}
\graphicspath{{figures/}}

\newtheorem{lemma}{Lemma}
\renewcommand{\thelemma}{S\arabic{lemma}}
\newcommand{\cvar}{\mathrm{CVaR}}
\newcommand{\E}{\mathbb{E}}
\newcommand{\Lt}{\widetilde L}
\newcommand{\Rt}{R_{\mathrm{tail}}}
\newcommand{\Rm}{R^{\mathrm{model}}}
\newcommand{\umean}{u_{\mathrm{mean}}}
\newcommand{\ucvar}{u_{\cvar}}
\newcommand{\figorbox}[2]{\includegraphics[width=#2]{#1}}
\newcommand{\pending}[1]{\textcolor{orange}{\textbf{[#1]}}}
\renewcommand{\thesection}{S\arabic{section}}
\renewcommand{\theequation}{S\arabic{equation}}
\renewcommand{\thefigure}{S\arabic{figure}}
\renewcommand{\thetable}{S\arabic{table}}
\makeatletter\renewcommand{\p@subsection}{\thesection.}\makeatother

\begin{document}
\title{Supplemental Material for\\ ``When does tail-risk optimization change a quantum gate?
An equivalence boundary for mean and CVaR control''}
\author{Sashreek Satapathy}
\date{\today}
\maketitle

Equations, figures, tables and sections of the main text are cited without
the prefix S (e.g.\ Eq.~(\ref{M-eq:ru}), Theorem~\ref{M-thm:main}); those of
this Supplement carry it. In the code and data files, experiments E1--E6 are
called Stage~0, E7--E8 Stage~1 and E9 Stage~2.

\tableofcontents
% =====================================================================
\section{Proofs}
\label{sec:proofs}

Throughout, $Z,Z_1,Z_2$ denote independent standard normal variables,
$\varphi$ and $\Phi$ the standard normal density and distribution function,
and $h(\phi)=\sin^2(\phi/2)$. We use four standard properties of
$\cvar_\alpha$, all immediate from the Rockafellar--Uryasev form
(\ref{M-eq:ru}): \emph{monotonicity} ($X\le Y$ almost surely implies
$\cvar_\alpha X\le\cvar_\alpha Y$), \emph{translation invariance}
($\cvar_\alpha(X+c)=\cvar_\alpha X+c$), \emph{positive homogeneity}
($\cvar_\alpha(sX)=s\,\cvar_\alpha X$ for $s\ge0$) and \emph{subadditivity}
($\cvar_\alpha(X+Y)\le\cvar_\alpha X+\cvar_\alpha Y$)
\cite{Rockafellar2000,Acerbi2002}. Subadditivity follows by choosing
$\tau=\tau_X+\tau_Y$ in (\ref{M-eq:ru}) and using
$(x+y-\tau_X-\tau_Y)_+\le(x-\tau_X)_++(y-\tau_Y)_+$. CVaR is also
\emph{law invariant}: it depends on $X$ only through its distribution.

\subsection{The equivalence theorem}

\paragraph*{Periodicity.} The infidelity $h$ is strictly increasing in $|\phi|$ on $[0,\pi]$, strictly decreasing on $[\pi,2\pi]$ and $2\pi$-periodic: $h'(\phi)=\tfrac12\sin\phi$, positive on
$(0,\pi)$ and negative on $(\pi,2\pi)$; $h$ is even and $2\pi$-periodic.
\qed

\paragraph*{Mean infidelity.} $h(\phi)=\tfrac12(1-\cos\phi)$ and, for
$\phi\sim\mathcal N(0,\sigma^2)$, $\E\cos\phi=\mathrm{Re}\,\E e^{i\phi}=e^{-\sigma^2/2}$
(the characteristic function of a Gaussian). Hence
$\E h(\phi)=\tfrac12(1-e^{-\sigma^2/2})$, whose derivative in $\sigma^2$ is
$\tfrac14e^{-\sigma^2/2}>0$. \qed

\paragraph*{Majorant.} $\Lt\ge h$ everywhere, with equality for $|\phi|\le\pi$, and $\Lt$ is non-decreasing in $|\phi|$. For $|\phi|\le\pi$, $\Lt=h$. For
$|\phi|>\pi$, $\Lt=1\ge h$. Monotonicity: on $[0,\pi]$, $\Lt=h$ increases
from 0 to 1 (periodicity, above); on $(\pi,\infty)$ it is constant at
its maximum value 1. \qed

\paragraph*{Sandwich bound, Eq.~(\ref{M-eq:sandwich}).} Let $E=\{|\phi|>\pi\}$ and $p=\Pr(E)$.
\emph{Upper bound}: $L\le\Lt$ pointwise , so
$\cvar_\alpha L\le\cvar_\alpha\Lt$ by monotonicity. \emph{Lower bound}:
pointwise, $\Lt\le L+\mathbf 1_E$, since on $E^c$ the two sides are equal
up to the nonnegative indicator, and on $E$ the left side is 1 while the right
side is at least 1. Monotonicity and subadditivity give
$\cvar_\alpha\Lt\le\cvar_\alpha L+\cvar_\alpha\mathbf 1_E$. For the Bernoulli
variable $\mathbf 1_E$, Eq.~(\ref{M-eq:ru}) gives
\[
\cvar_\alpha\mathbf 1_E=\min_\tau\Big\{\tau+\tfrac1\alpha\big[p(1-\tau)_++(1-p)(-\tau)_+\big]\Big\}
=\min\!\Big(1,\frac p\alpha\Big)\le\frac p\alpha ,
\]
where the minimum is attained at $\tau=0$ if $p\le\alpha$ and at $\tau=1$
otherwise (the objective is piecewise linear in $\tau$ with kinks at 0 and 1).
Rearranging,
$\cvar_\alpha L\ge\cvar_\alpha\Lt-p/\alpha$. No condition on $p$ versus
$\alpha$ is needed, because the bound $\min(1,p/\alpha)\le p/\alpha$ holds
always. \qed

\paragraph*{Theorem~\ref{M-thm:main}.} Write $\phi_i=\sigma_{f_i}Z$ (same $Z$;
allowed because CVaR is law invariant). By the lower half of
the sandwich bound for $f_2$ and the upper half for $f_1$,
\[
\cvar_\alpha L(\phi_2)-\cvar_\alpha L(\phi_1)\ge
\big[\cvar_\alpha\Lt(\phi_2)-r_\alpha(\sigma_{f_2})\big]-\cvar_\alpha\Lt(\phi_1)
=\Delta_{\mathrm{mono}}-r_\alpha(\sigma_{f_2}).
\]
$\Delta_{\mathrm{mono}}\ge0$: since $\sigma_{f_2}>\sigma_{f_1}$,
$|\phi_2|\ge|\phi_1|$ for every realization of $Z$, and $\Lt$ is
non-decreasing in $|\phi|$, so $\Lt(\phi_2)\ge\Lt(\phi_1)$ pointwise; apply
monotonicity. For the residual, $\Pr(|\phi|>\pi)=2\Pr(Z>\pi/\sigma)$ and the
Mills-ratio bound $\Pr(Z>x)\le\varphi(x)/x$ for $x>0$ (which follows from
$\int_x^\infty\varphi(t)\,dt\le\int_x^\infty\tfrac tx\varphi(t)\,dt=\varphi(x)/x$)
gives
\[
r_\alpha(\sigma)=\frac{2\Pr(Z>\pi/\sigma)}{\alpha}\le\frac{2}{\alpha}\,\frac{\sigma}{\pi}\,\frac{e^{-\pi^2/(2\sigma^2)}}{\sqrt{2\pi}}
=\frac1\alpha\sqrt{\frac2\pi}\,\frac\sigma\pi\,e^{-\pi^2/(2\sigma^2)} .\qquad\qed
\]

\paragraph*{When is $\Delta_{\mathrm{mono}}>0$ strictly?} If $\sigma_{f_1}<\sigma_{f_2}$ and
$\Pr(|\phi_2|<\pi)>0$ (always true for a Gaussian), then
$\Lt(\phi_2)>\Lt(\phi_1)$ on the event $\{0<|\phi_2|<\pi\}$, which has
positive probability and includes part of the upper $\alpha$ tail of
$\Lt(\phi_2)$ whenever $\Pr(|\phi_2|>\pi)<\alpha$. The CVaR difference is then
strictly positive. In the small-noise regime it equals
$c_\alpha(\sigma_{f_2}^2-\sigma_{f_1}^2)+O(\sigma^4)$ by
Proposition~\ref{M-prop:ratio}.

\paragraph*{Optimizer equivalence.} \emph{Claim:} let $u^\star$ minimize $\E L(\phi[u])$. If $\Delta_{\mathrm{mono}}(\sigma_{u^\star},\sigma_u)>r_\alpha(\sigma_u)$ for every control with $\sigma_u>\sigma_{u^\star}$, then $u^\star$ is a global CVaR minimizer, and every CVaR minimizer has the same loss distribution. \emph{Proof:} By the mean formula, $\E L(\phi[u])$ is a
strictly increasing function of $\sigma_u$, so $u^\star$ minimizes $\sigma_u$.
Take any control $u$. If $\sigma_u>\sigma_{u^\star}$, Theorem~\ref{M-thm:main}
and the hypothesis give
$\cvar_\alpha L(u)-\cvar_\alpha L(u^\star)\ge\Delta_{\mathrm{mono}}-r_\alpha>0$.
If $\sigma_u=\sigma_{u^\star}$, the two losses have the same distribution and
the same CVaR. Hence $u^\star$ is a global CVaR minimizer, and every CVaR
minimizer has $\sigma_u=\sigma_{u^\star}$. \qed

\paragraph*{Size of the residual.} At $\sigma^2=0.05$,
$\pi/\sigma=14.05$ and $\pi^2/(2\sigma^2)=98.7$, so
$2\Pr(Z>14.05)\approx7.7\times10^{-45}$ and
$r_{0.05}\approx1.5\times10^{-43}$.

\paragraph*{Proposition~\ref{M-prop:ratio}.} $h(\phi)=\phi^2/4-\phi^4/48+\dots$,
so $|h(\sigma Z)-\sigma^2Z^2/4|\le\sigma^4Z^4/48$. CVaR is 1-Lipschitz in the
sense $|\cvar_\alpha X-\cvar_\alpha Y|\le\cvar_\alpha|X-Y|$ (monotonicity plus
subadditivity), and $\cvar_\alpha(Z^4)<\infty$, so
$\cvar_\alpha h(\sigma Z)=\tfrac{\sigma^2}{4}\cvar_\alpha(Z^2)+O(\sigma^4)$. The
upper $\alpha$ tail of $Z^2$ is the event $\{|Z|>z_\alpha\}$,
$z_\alpha=\Phi^{-1}(1-\alpha/2)$, of probability exactly $\alpha$, so
\[
\cvar_\alpha(Z^2)=\frac1\alpha\,2\!\int_{z_\alpha}^\infty t^2\varphi(t)\,dt
=\frac2\alpha\Big[z_\alpha\varphi(z_\alpha)+\big(1-\Phi(z_\alpha)\big)\Big]
=1+\frac{2z_\alpha\varphi(z_\alpha)}{\alpha},
\]
since $-t\varphi(t)+\Phi(t)$ is an antiderivative of $t^2\varphi(t)$
(integrate by parts with $\varphi'=-t\varphi$).
Numerically, for $\alpha=0.20,0.10,0.05,0.025,0.01$:
$z_\alpha=1.282,1.645,1.960,2.241,2.576$ and
$\cvar_\alpha(Z^2)=3.25,4.39,5.58,6.80,8.45$. The mean has $\E Z^2=1$, hence
the ratios. \qed

\paragraph*{Average gate fidelity.} For a unitary error $U_{\mathrm{err}}$
on a $d$-level system, $F_{\mathrm{avg}}=(dF_e+1)/(d+1)$ with entanglement
fidelity $F_e=|\mathrm{Tr}\,U_{\mathrm{err}}/d|^2$ \cite{Nielsen2002}. For
$U_{\mathrm{err}}=e^{-i\phi Z/2}$ and $d=2$, $F_e=\cos^2(\phi/2)$, so
$1-F_{\mathrm{avg}}=\tfrac23(1-\cos^2(\phi/2))=\tfrac23h(\phi)$. Taking
expectations with the mean formula gives
$\tfrac13(1-e^{-\sigma_f^2/2})$. Since $\tfrac23h$ is a positive multiple of
$h$, every step of Theorem~\ref{M-thm:main} goes through unchanged. \qed

\label{sec:proofs-ext}\paragraph*{Temporal correlation.} \emph{Claim:} for any stationary Gaussian noise, e.g.\ Ornstein--Uhlenbeck noise $C(t,t')=\sigma^2e^{-|t-t'|/\tau_c}$, $\phi[u]$ is exactly Gaussian and Theorem~\ref{M-thm:main} applies verbatim. \emph{Proof:} $\phi[f]=\int f\eta\,dt$ is a
continuous linear functional of a Gaussian process, hence Gaussian, with
variance $\int\!\!\int f(t)C(t,t')f(t')\,dt\,dt'$. For $f\equiv1$ and
$C=\sigma^2e^{-|t-t'|/\tau_c}$,
\[
\sigma_f^2=\sigma^2\!\int_0^{T_g}\!\!\int_0^{T_g}e^{-|t-t'|/\tau_c}dt\,dt'
=2\sigma^2\!\int_0^{T_g}(T_g-s)e^{-s/\tau_c}ds
=2\sigma^2\tau_c\big(T_g-\tau_c(1-e^{-T_g/\tau_c})\big).
\]
As $\tau_c\to\infty$ this tends to $\sigma^2T_g^2$, the quasi-static limit
$\sigma_\phi^2$ used in this paper. \qed

\paragraph*{Several channels on one axis.} If $\boldsymbol\eta(t)$ is a jointly
Gaussian vector process and $\phi=\sum_j\int f_j\eta_j\,dt$, then $\phi$ is a
linear functional of a Gaussian process, hence Gaussian, with variance
\[Q(f)=\sum_{jk}\int\!\!\int f_j(t)C_{jk}(t,t')f_k(t')\,dt\,dt' .\qquad\qed\]

\paragraph*{Corollary~\ref{M-cor:general}.} Let $\tau(u_1)\le\tau(u_2)$ and
$L_{u_i}=h(\tau(u_i)Z)$ with $h$ non-decreasing in $|\cdot|$. Then
$L_{u_1}\le L_{u_2}$ pointwise on a common probability space, which is
equivalent to $L_{u_1}\le_{\mathrm{st}}L_{u_2}$ \cite{Shaked2007}. Any
risk measure that is monotone and law invariant therefore satisfies
$\varrho(L_{u_1})\le\varrho(L_{u_2})$, so all such measures are minimized by
the controls with the smallest $\tau(u)$. \qed

\subsection{Two non-parallel channels}
\label{sec:proofs-twoaxis}

\paragraph*{Reduction to two weights.} To leading order the error of a draw is
$\boldsymbol\epsilon=\delta\,\mathbf b+\gamma\,\mathbf a$ and the loss is
$\boldsymbol\epsilon^{\mathsf T}W\boldsymbol\epsilon=\mathbf x^{\mathsf T}M\mathbf x$
with $\mathbf x=(\delta/\sigma_\delta,\gamma/\sigma_\gamma)$ and
$M=G^{\mathsf T}WG$, $G=(\sigma_\delta\mathbf b,\sigma_\gamma\mathbf a)$. With
$\mathbf x=C^{1/2}\mathbf y$, $\mathbf y$ standard normal and $C$ the
correlation matrix, $L=\mathbf y^{\mathsf T}C^{1/2}MC^{1/2}\mathbf y$.
Diagonalizing the symmetric positive semidefinite matrix $C^{1/2}MC^{1/2}$ by
an orthogonal matrix (which maps $\mathbf y$ to another standard normal
vector) gives $L=\lambda_1Z_1^2+\lambda_2Z_2^2$.

\paragraph*{Proposition~\ref{M-prop:twoaxis}(ii), Schur-convexity.} Fix $\tau$
and write $f_\tau(\lambda)=\tau+\alpha^{-1}\E(\lambda_1X_1+\lambda_2X_2-\tau)_+$
with $X_i=Z_i^2$. For each outcome, $\lambda\mapsto(\lambda\cdot X-\tau)_+$ is
convex (a convex function of an affine function), and expectation preserves
convexity. Since $X_1,X_2$ are i.i.d., $f_\tau(\lambda_1,\lambda_2)=f_\tau(\lambda_2,\lambda_1)$.
A symmetric convex function is Schur-convex \cite[Ch.~3]{Marshall2011}: if
$\lambda$ majorizes $\mu$ (same sum, $\mu$ more equal), then $\mu$ is a convex
combination of $\lambda$ and its permutation, $\mu=t\lambda+(1-t)\lambda^\pi$,
so $f_\tau(\mu)\le tf_\tau(\lambda)+(1-t)f_\tau(\lambda^\pi)=f_\tau(\lambda)$.
This holds for every $\tau$, so taking the minimum over $\tau$ on both sides
gives $\cvar_\alpha(L_\mu)\le\cvar_\alpha(L_\lambda)$. \qed

\paragraph*{The two extreme values.} \emph{Equal split}, $\lambda=(\tfrac12,\tfrac12)$:
$L=\tfrac12(Z_1^2+Z_2^2)$ is exponential with unit mean. By memorylessness,
$\E[L\,|\,L>q]=q+1$, and $\Pr(L>q)=\alpha$ at $q=\ln(1/\alpha)$, so
$\cvar_\alpha=1+\ln(1/\alpha)$, which is $4.00$ at $\alpha=0.05$.
\emph{One axis}, $\lambda=(1,0)$: $L=Z^2$ and, from the proof of
Proposition~\ref{M-prop:ratio},
$\cvar_\alpha=1+2c\varphi(c)/\alpha=2[c\varphi(c)+\alpha/2]/\alpha$ with
$c=z_\alpha$, which is $5.58$ at $\alpha=0.05$. The difference is
$1-4.00/5.58=28\%$ of the one-axis value.
Figure~\ref{M-fig:theory}(c) shows the full curve for four values of $\alpha$,
computed by Monte Carlo with $6\times10^5$ draws.

\paragraph*{Proposition~\ref{M-prop:twoaxis}(iii).} Suppose
$L_u\ge_{\mathrm{st}}L_v$ with $\E L_u=\E L_v$. Realize them with
$L_u\ge L_v$ pointwise; then $L_u-L_v\ge0$ has zero mean, so $L_u=L_v$ almost
surely, and the two laws coincide. For two weight pairs with equal sum and
different values, $\mathrm{Var}\,L=2(\lambda_1^2+\lambda_2^2)$ differs (since
$\mathrm{Var}\,Z^2=2$), so the laws differ and neither dominates.
\emph{Sufficient condition.} If the sorted weights satisfy
$\lambda_{(1)}\le\mu_{(1)}$ and $\lambda_{(2)}\le\mu_{(2)}$, then (relabelling
the independent $Z_i$) $\lambda_{(1)}Z_1^2+\lambda_{(2)}Z_2^2\le\mu_{(1)}Z_1^2+\mu_{(2)}Z_2^2$
pointwise, hence dominance. A family with a common ratio
$\lambda_2/\lambda_1=r$ is $\lambda=s(1,r)$, totally ordered by $s$, so it
satisfies the condition. \emph{Not necessary.} The pairs $(2,1)$ and
$(1,0.9)$ have different ratios ($0.5$ and $0.9$) but are componentwise
ordered, so $(1,0.9)$ is dominated by $(2,1)$. An earlier draft of this paper
stated that a common ratio was also necessary; this example shows it is not.
\qed

\subsection{The effective dimension}
\label{sec:proofs-meff}

\paragraph*{Parts (i) and (ii).} Let $X\sim\chi^2_m$ and $q=\chi^2_{m,1-\alpha}$.
The size-biasing identity for the chi-square distribution,
\begin{equation}
\E\big[X\,\mathbf 1\{X>q\}\big]=m\,P\big(\chi^2_{m+2}>q\big),
\label{eq:sizebias}
\end{equation}
follows from $x\,f_m(x)=m\,f_{m+2}(x)$ for the densities $f_m$, which is
immediate from $f_m(x)=x^{m/2-1}e^{-x/2}/(2^{m/2}\Gamma(m/2))$ and
$\Gamma(m/2+1)=(m/2)\Gamma(m/2)$. Since $\E X=m$ and $X$ is continuous,
\begin{equation}
\frac{\cvar_\alpha(X)}{\E X}=\frac{\E[X\mathbf 1\{X>q\}]}{\alpha\,\E X}
=\frac1\alpha P\big(\chi^2_{m+2}>q\big)=\rho_\alpha(m),
\end{equation}
and the ratio is invariant under the rescaling $L=(\sum_iw_i/m)X$ that turns
equal weights into $\chi^2_m$, giving (i).

For $m=1$, $q=z_\alpha^2$ with $z_\alpha=\Phi^{-1}(1-\alpha/2)$, and
$P(\chi^2_3>z_\alpha^2)=\alpha+2z_\alpha\varphi(z_\alpha)$ by integrating the
$\chi^2_3$ density by parts; dividing by $\alpha$ gives
$\rho_\alpha(1)=1+2z_\alpha\varphi(z_\alpha)/\alpha=r_\alpha$, i.e.\
Proposition~\ref{M-prop:ratio}. For $m=2$, $X/2$ is exponential with unit
mean, $q/2=\ln(1/\alpha)$, and the memorylessness of the exponential gives
$\E[X/2\mid X/2>\ln(1/\alpha)]=1+\ln(1/\alpha)$, i.e.\ $\rho_\alpha(2)$ is the
equal-split value of Proposition~\ref{M-prop:twoaxis}. The two propositions of
the main text are thus the $m=1$ and $m=2$ points of one curve.

\paragraph*{Part (iii).} $\chi^2_m/m$ has mean $1$ and variance $2/m$ and
concentrates at $1$, so its $\alpha$-tail average tends to $1$; strict
monotonicity in $m$ follows because $\chi^2_{m}/m$ for $m'>m$ is a strict
mean-preserving contraction of $\chi^2_{m'}/m'$ in the convex order, and
$\cvar_\alpha$ is strictly monotone under that order for a continuous
distribution \cite{Shaked2007}. Equation~(\ref{M-eq:rhoeff}) extends to
non-integer $m$ by the same formula, and is smooth and decreasing there.

\paragraph*{Part (iv): the participation-ratio approximation.} For unequal
weights we replace the spectrum $\{w_i\}$ by the equal-weight spectrum with
the same first and second moments, which has exactly $m_{\mathrm{eff}}$
elements by Eq.~(\ref{M-eq:meff}); $L$ and its surrogate then share their mean
and variance, and the approximation amounts to matching the first two
cumulants of the loss. It is not exact, because the tail of a weighted
chi-square sum is eventually governed by the largest weight alone, so the
error must grow as $\alpha\to0$. We quantify it on $400$ random spectra of
$2$--$10$ directions drawn in four styles (uniform weights; uniform cubed;
log-uniform over four decades; and one dominant weight with small companions),
each evaluated with $4\times10^5$ draws.

\begin{table}[h]
\caption{Relative error of $\rho_\alpha(m_{\mathrm{eff}})$ against the true
ratio over $400$ random spectra. Negative means the approximation
underestimates the ratio, i.e.\ is conservative for a tail objective.}
\label{tab:prop5acc}
\begin{ruledtabular}
\begin{tabular}{lccccc}
$\alpha$ & 0.20 & 0.10 & 0.05 & 0.025 & 0.01\\
\hline
median error (\%) & 0.42 & -0.80 & -1.94 & -3.01 & -4.24\\
95th pct.\ $|$error$|$ (\%) & 1.16 & 2.28 & 5.29 & 7.78 & 10.37\\
worst error (\%) & 1.30 & -3.07 & -6.51 & -9.53 & -12.80\\
fraction conservative & 0.22 & 0.97 & 0.99 & 0.99 & 0.99\\
\end{tabular}
\end{ruledtabular}
\end{table}

At $\alpha=0.05$ the approximation is conservative at $99\%$ of spectra and
within $6.5\%$ everywhere; at $\alpha=0.01$ the worst case reaches $12.8\%$.
For the controls of Sec.~\ref{sec:stage4}, where $m_{\mathrm{eff}}\le3$ and the
spectra are mild, the observed error is a median $2.6\%$. Part (iv) is
therefore an empirical statement with a stated domain, not a theorem, and the
main text marks it as such.

\subsection{The rare-channel mixture}
\label{sec:proofs-mixture}

\paragraph*{Proposition~\ref{M-prop:mixture}, with the error term.} Let
$L=L_c+\ell_j\mathbf 1_J$ with $J$ independent of $L_c$, $P(J)=p<\alpha$, and
let $v_\gamma$ denote the $(1-\gamma)$-quantile of $L_c$. Assume
$\ell_j>v_{\alpha-p}$, so that every point of $J$ lies above every point of
$J^c$ that enters the tail.

Write $\gamma=\alpha-p$. The worst $\alpha$ of the mass of $L$ is then exactly
$J$ together with the worst $\gamma$ of $J^c$, and since $\cvar$ at level
$\alpha$ is the average of the worst $\alpha$ fraction,
\begin{equation}
\alpha\,\cvar_\alpha(L)=p\big(\ell_j+\E[L_c\mid J]\big)
+\gamma\,\cvar_{\gamma}(L_c\mid J^c).
\label{eq:mixexact}
\end{equation}
Independence gives $\E[L_c\mid J]=\E L_c$ and
$\cvar_\gamma(L_c\mid J^c)=\cvar_\gamma(L_c)$, so with
$\beta=p\ell_j/\E L_c$ and $\E L=(1+\beta)\E L_c$,
\begin{equation}
\frac{\cvar_\alpha(L)}{\E L}
=\frac{1}{1+\beta}\left[\frac{\beta}{\alpha}+\frac{p}{\alpha}
+\frac{\alpha-p}{\alpha}\,\frac{\cvar_{\alpha-p}(L_c)}{\E L_c}\right],
\label{eq:mixexact2}
\end{equation}
which is exact. Equation~(\ref{M-eq:mixture}) is Eq.~\eqref{eq:mixexact2} with
$p/\alpha\to0$. Two corrections are dropped. The term $p/\alpha$ is the jump
branch's share of the \emph{continuous} loss, a relative error of at most
$(p/\alpha)/\rho_\alpha$. The second is the difference between
$\frac{\alpha-p}{\alpha}\,\cvar_{\alpha-p}(L_c)/\E L_c$ and $\rho_\alpha$;
since $\gamma\,\cvar_\gamma$ is the tail integral, which is Lipschitz in
$\gamma$ with constant $v_\gamma$, that difference is at most
$p\,v_{\alpha-p}/(\alpha\,\E L_c)$. Both are $O(p/\alpha)$ and both are
positive, so Eq.~(\ref{M-eq:mixture}) \emph{underestimates} the ratio.
Throughout Sec.~\ref{sec:stage5}, $p\le2.0\times10^{-3}$ at $\alpha=0.05$, so
$p/\alpha\le0.04$ and the bound on the relative error stays below $1.5\%$ at
every point --- consistent with the observed median of $1.2\%$ and with its
sign.

\paragraph*{Two limits and one warning.} As $\beta\to0$ the ratio returns to
$\rho_\alpha(m_{\mathrm{eff}})$. As $\beta\to\infty$ it approaches $1/\alpha$,
the largest ratio any loss can have when its whole tail mass sits in a single
atom, so the mixture law interpolates between the continuous-noise prediction
and that ceiling. The condition $p<\alpha$ matters: for $p>\alpha$ the worst
$\alpha$ of the mass lies entirely inside $J$, the continuous part drops out
of the tail, and the ratio saturates at $\ell_j/\E L$ instead. The evaluation
grid of Sec.~\ref{sec:stage5} deliberately spans
$\alpha\in\{0.5,\,0.2,\,0.05,\,0.01,\,2\times10^{-3},\,5\times10^{-4},\,
10^{-4}\}$ so that both sides of $p\approx\alpha$ are covered.

\subsection{The reduced second-order model}

\paragraph*{Proposition~\ref{M-prop:reduced}.}
\begin{proof}
(i) $\E L$ is quadratic in $\mathbf e_0$ with Hessian $2W$ and stationary at
$\mathbf e_0=-\E[A\mathbf z+Q(\mathbf z)]=-\E\,Q(\mathbf z)$.
(ii) $\boldsymbol\epsilon=\|A\|\,(\hat A\mathbf z+\kappa\hat Q(\mathbf z)+\mathbf e_0/\|A\|)$;
the factor $\|A\|^2$ cancels in the ratio defining $\Rm_\alpha$, and
$\mathbf e_0/\|A\|$ ranges over $\mathbb R^3$ as $\mathbf e_0$ does. Scaling
the noise by $s$ sends $A\to sA$ and $Q\to s^2Q$.
(iii) At $\kappa=0$ the noise part $\hat A\mathbf z$ is symmetric under
$\mathbf z\to-\mathbf z$, so $\cvar_\alpha(L(\mathbf e_0))$ is even in
$\mathbf e_0$. It is also convex, since $L$ is convex in $\mathbf e_0$ for
every draw and CVaR is monotone and convex. Its minimum is therefore at
$\mathbf e_0=0$, which is also the mean optimum.
\end{proof}

Two additional remarks:

\emph{(a) Why $\mathbf e_0/\|A\|$ is free.} In (ii), the minimization of
each objective runs over all $\mathbf e_0\in\mathbb R^3$, and
$\mathbf e_0\mapsto\mathbf e_0/\|A\|$ is a bijection of $\mathbb R^3$, so the
optimal \emph{normalized} offset, and therefore both optimal values up to the
common factor $\|A\|^2$, depend only on $(\kappa,\hat A,\hat Q)$.

\emph{(b) What (ii) does not say.} The scaling $\kappa\mapsto s\kappa$ holds
for a \emph{fixed} $A$ and $Q$, i.e.\ a fixed control. If the control is
re-optimized at each noise level, $A$ and $Q$ change as well. The prediction
(c) of the prospective test (Sec.~\ref{M-sec:mechanism}) overlooked this, and
failed for that reason.

\subsection{Remark~\ref{M-rem:precomp}: the one-dimensional pre-compensation model}
\label{sec:proofs-precomp}

Let $\boldsymbol\epsilon=\mathbf e_0+\mathbf qz^2$ and restrict the offset to
the line $\mathbf e_0=-c\,\mathbf q$ (components of $\mathbf e_0$
perpendicular to $\mathbf q$ only add a constant $|\mathbf e_\perp|^2_W\ge0$
to every draw and are optimally zero for both objectives, by translation
invariance). With the loss metric $W$, $L=w\,(z^2-c)^2$ with
$w=\mathbf q^{\mathsf T}W\mathbf q>0$, and $w$ drops out of both argmins.

\emph{Mean.} $\E(z^2-c)^2=\mathrm{Var}(z^2)+(\E z^2-c)^2=2+(1-c)^2$, minimized
at $c=1$ with value 2.

\emph{CVaR.} By Eq.~(\ref{M-eq:ru}), $\cvar_\alpha L(c)=\min_\tau g(c,\tau)$,
and $g$ is jointly convex in $(c,\tau)$, so the optimum satisfies the
first-order conditions. Differentiating under the expectation,
$\partial_cg=\alpha^{-1}\E[\partial_cL\;\mathbf 1\{L>\tau\}]=-2\alpha^{-1}\E[(z^2-c)\mathbf 1\{L>\tau\}]$,
and at the optimal $\tau$ (the VaR) the condition $\partial_cg=0$ reads
\begin{equation}
\E\big[z^2-c\;\big|\;\text{tail}\big]=0 :
\label{eq:precomp-foc}
\end{equation}
the \emph{average signed error over the tail draws vanishes}. The mean
condition is the same with the whole distribution in place of the tail.
Solving numerically ($2\times10^6$ draws) gives $c=2.76$ at $\alpha=0.05$
($\cvar=13.8$, against $24.2$ at $c=1$) and $c=1.70$ at $\alpha=0.2$. The
structure of the solution is instructive. At the optimum the VaR equals
almost exactly $c^2$, the loss of a draw with $z=0$, so the tail contains two
kinds of draws: the \emph{largest} draws, with $z^2>2c$, whose error is
under-compensated (1.9\% of all draws at $\alpha=0.05$), and the
\emph{smallest} draws, with $z\approx0$, whose error is over-compensated (the
remaining 3.1\%; 62\% of the tail). Condition \eqref{eq:precomp-foc} balances
the two. The CVaR optimum is therefore the offset at which the worst
over-compensated and the worst under-compensated draws are equally bad, which
is why it lies much further out than the typical draw.


\subsection{The $\kappa^2$ law}
\label{sec:proofs-kappa2}

Throughout this section the reduced model of
Proposition~\ref{M-prop:reduced} is used in its normalized form: with
$a=\|A\|_F$, $\hat A=A/a$ and $\hat Q=Q/(a\kappa)$,
\[
\bm\epsilon=a\big(\hat A\mathbf z+\kappa\hat Q(\mathbf z)+\hat{\mathbf e}_0\big),
\qquad L=\bm\epsilon^{\mathsf T}W\bm\epsilon ,
\]
and the factor $a^2$ cancels in $R_\alpha$. We drop the hats and write the
offset as $\mathbf e_0=\kappa\mathbf f$, which is the scale at which both
optima live: the mean optimum is exactly
$\mathbf e_0=-\kappa\,\E\,Q(\mathbf z)$.

\paragraph*{Expansion of the loss.} With $L_0=\mathbf z^{\mathsf T}M\mathbf z$,
$M=A^{\mathsf T}WA$,
\begin{equation}
L=L_0+\kappa\,h_{\mathbf f}(\mathbf z)+\kappa^2g_{\mathbf f}(\mathbf z),\qquad
\begin{aligned}
h_{\mathbf f}&=2(A\mathbf z)^{\mathsf T}W\big(Q(\mathbf z)+\mathbf f\big),\\
g_{\mathbf f}&=\big(Q+\mathbf f\big)^{\mathsf T}W\big(Q+\mathbf f\big).
\end{aligned}
\label{eq:Lexp}
\end{equation}
$L_0$ and $g_{\mathbf f}$ are even in $\mathbf z$; $h_{\mathbf f}$ is odd,
being a product of a linear and a quadratic form.

\paragraph*{First order vanishes.} Let $\tau^\ast(\kappa)$ attain the minimum
in Eq.~(\ref{M-eq:ru}). Since the objective
$\tau+\alpha^{-1}\E(L-\tau)_+$ is jointly convex and differentiable in
$\kappa$ for fixed $\tau$, the envelope theorem gives
\begin{equation}
\frac{d}{d\kappa}\cvar_\alpha(L)=\frac1\alpha\,\E\big[(h_{\mathbf f}+2\kappa g_{\mathbf f})\,\mathbf 1\{L>\tau^\ast\}\big].
\label{eq:env}
\end{equation}
At $\kappa=0$ the tail event is $\{L_0>\tau_0\}$, which is invariant under
$\mathbf z\to-\mathbf z$, while $h_{\mathbf f}$ is odd; the expectation
vanishes. \emph{The $\alpha$-tail of a symmetric loss has two sides, and to
first order a coherent offset gains on one exactly what it loses on the
other.} This cancellation, and not any smallness of the offset itself, is why
a threshold in $\kappa$ exists.

\paragraph*{Second order.} Differentiating Eq.~\eqref{eq:env} once more at
$\kappa=0$, using the value-at-risk sensitivity
$d\tau^\ast/d\kappa=\E[h_{\mathbf f}\mid L_0=\tau_0]$ (valid when $L_0$ has a
continuous positive density $\varrho$ at $\tau_0$, which holds for a
nondegenerate Gaussian quadratic form), the level-set term contributes the
conditional variance of $h_{\mathbf f}$ and
\begin{equation}
\frac{d^2}{d\kappa^2}\cvar_\alpha(L)\Big|_{0}
=\frac2\alpha\E\big[g_{\mathbf f}\mathbf 1\{L_0>\tau_0\}\big]
+\frac{\varrho(\tau_0)}{\alpha}\mathrm{Var}\big[h_{\mathbf f}\mid L_0=\tau_0\big],
\end{equation}
which is Eq.~(\ref{M-eq:Jf}) after division by two. Both terms are quadratic
in $\mathbf f$, with Hessian
\begin{equation}
H=2W+\frac{4\varrho(\tau_0)}{\alpha}\,\Sigma_u,\qquad
\Sigma_u=\mathrm{Cov}\big[WA\mathbf z\mid L_0=\tau_0\big]\succeq0,
\label{eq:Hess}
\end{equation}
so $J$ is strictly convex ($W\succ0$) and its minimizer is explicit,
$\mathbf f^\star=-H^{-1}\mathbf b$ with
$\mathbf b=2W\,\E[Q\mid L_0>\tau_0]\,\alpha/\alpha+(2\varrho/\alpha)\mathrm{Cov}[h_0,WA\mathbf z\mid L_0=\tau_0]$
and $h_0=h_{\mathbf 0}$. Since $J$ is quadratic,
\begin{equation}
J(\mathbf f_{\mathrm{mean}})-\min_{\mathbf f}J
=\tfrac12(\mathbf f_{\mathrm{mean}}-\mathbf f^\star)^{\mathsf T}H(\mathbf f_{\mathrm{mean}}-\mathbf f^\star),
\end{equation}
and dividing by $\cvar_\alpha(L_0)$ gives Eq.~(\ref{M-eq:kappa2}). Both optima
are $O(\kappa)$, so the remainder is $O(\kappa^3)$.

\paragraph*{One dominant direction.} Let $n=1$ and let the second-order error
be parallel to the first-order one, $\bm\epsilon=(e_0+z+\kappa z^2)\mathbf d$.
Then $W$ enters only through $\mathbf d^{\mathsf T}W\mathbf d>0$, which cancels
in $R_\alpha$, and with $e_0=-\kappa c$ the loss is $L=F(z)^2$,
$F(z)=z+\kappa(z^2-c)$. Write $s=z_\alpha$, so that
$2\big(1-\Phi(s)\big)=\alpha$. The tail of $L$ is
$\{z>z_+\}\cup\{z<z_-\}$ with $F(z_\pm)=\pm V$, and expanding
$z_+=s+\kappa a_1+\kappa^2a_2$, $z_-=-s+\kappa b_1+\kappa^2b_2$,
$V=s+\kappa v_1+\kappa^2v_2$ in Eq.~\eqref{eq:Lexp} gives
\begin{equation*}
a_1+s^2-c=v_1,\qquad b_1+s^2-c=-v_1,
\end{equation*}
while the tail probability constraint
$\big[1-\Phi(z_+)\big]+\Phi(z_-)=\alpha$ reads, to first order,
$-a_1+b_1=0$. Hence $v_1=0$ and
\begin{equation}
a_1=b_1=c-s^2 :
\label{eq:a1}
\end{equation}
the two ends of the tail move by the same amount and in the same direction.
At second order the same constraint gives $b_2-a_2=-s\,a_1^2$. Collecting the
three contributions to $\alpha\cvar_\alpha(L)$ --- the integral of $F^2\varphi$
over the symmetric region $|z|>s$, the boundary shifts, and the curvature of
$F^2\varphi$ at $\pm s$ --- the terms in $s^2$ cancel and
\begin{equation}
\alpha\,\cvar_\alpha(L)=\underbrace{\E\big[z^2\mathbf 1_{|z|>s}\big]}_{=\;2s\varphi(s)+\alpha\;=\;\alpha r_\alpha}
+\kappa^2\Big[m_2(c)+2s\varphi(s)\,(c-s^2)^2\Big]+O(\kappa^3),
\label{eq:scalar}
\end{equation}
with $m_2(c)=\E[(z^2-c)^2\mathbf 1_{|z|>s}]$. The bracket is a quadratic in
$c$ with leading coefficient $\alpha+2s\varphi(s)=\alpha r_\alpha$ and
stationary point
\begin{equation*}
c^\star=\frac{\E[z^2\mathbf 1_{|z|>s}]+2s^3\varphi(s)}{\alpha+2s\varphi(s)}
=1+s^2\Big(1-\frac1{r_\alpha}\Big),
\end{equation*}
using $2s\varphi(s)=\alpha(r_\alpha-1)$. Therefore
$\cvar_\alpha(L)=r_\alpha\big[1+\kappa^2(c-c^\star)^2\big]+\text{const}+O(\kappa^3)$,
and since the mean-optimal offset is $c=1$,
\begin{equation}
R_\alpha(\kappa)=\frac{\cvar_\alpha(L)|_{c=1}-\cvar_\alpha(L)|_{c=c^\star}}{\cvar_\alpha(L)|_{c=1}}
=(c^\star-1)^2\kappa^2+O(\kappa^3),
\end{equation}
the factor $r_\alpha$ cancelling. This is Eq.~(\ref{M-eq:cstar}).

\paragraph*{The two limits of the offset.} $c^\star_\alpha\to1+z_\alpha^2(1-1/r_\alpha)$
as $\kappa\to0$ and $c\to2.76$ ($\alpha=0.05$) or $1.70$ ($\alpha=0.2$) as
$\kappa\to\infty$ (Remark~\ref{M-rem:precomp}), so the optimal
pre-compensation \emph{decreases} with $\kappa$: $4.15\to2.76$ at
$\alpha=0.05$. Direct minimization of the scalar model reproduces this,
$c^\star=4.20,\,4.17,\,3.82,\,3.16,\,2.86$ at
$\kappa=0.02,\,0.1,\,0.5,\,1.2,\,2.2$.

\paragraph*{Numerical verification.} \texttt{make\_kappa\_bound.py}
(Sec.~\ref{sec:repro}) checks every step without the simulator: (i) in the scalar
model the measured $R_\alpha/\kappa^2$ is $9.69$ and $1.29$ at $\kappa=0.05$
against the predicted $9.94$ and $1.29$ for $\alpha=0.05$ and $0.2$, and the
measured $c^\star$ is $4.14$ and $2.14$ against $4.15$ and $2.14$;
(ii) for random three-dimensional geometries the quadratic form
$J(\mathbf f)$ of Eq.~(\ref{M-eq:Jf}) agrees with direct Monte Carlo
evaluation of $[\cvar_\alpha(\kappa,\mathbf f)-\cvar_\alpha(0)]/\kappa^2$ to
about $1\%$ at $\kappa=0.02$--$0.04$; (iii) applied to the fitted $A$ and $Q$
of every diagnosed mean control, the closed form reproduces the numerically
optimized reduced model with median ratio $1.07$ (interquartile range
$0.94$--$1.12$) over the 30 controls with $\kappa<0.25$
[Fig.~\ref{M-fig:kappabound}(b)].

% =====================================================================
\section{Model, simulators and optimization}
\label{sec:modeldetails}
\begin{table}[htbp]
\centering\small
\begin{tabular}{@{}llp{0.42\textwidth}@{}}
\toprule
Quantity & Value & Note\\
\midrule
transmon levels $N$ & 4 (or 2) & 2 levels: leakage ablation\\
anharmonicity $\alpha_{\mathrm{anh}}/2\pi$ & $-200$~MHz & typical fixed-frequency transmon \cite{Koch2007,Krantz2019}\\
qubit frequency & 5~GHz & enters only the thermal population\\
temperature & 50~mK & thermal occupation $n_{\mathrm{th}}\approx0.0083$\\
$T_1$, $T_2^*$ & $45.5~\mu$s, $16.8~\mu$s & $\gamma_\varphi=1/T_2^*-1/(2T_1)$\\
target gate & $X_{\pi/2}=(I-iX)/\sqrt2$ & on the computational subspace\\
gate time $T_g$ & 6.3--32.5~ns & baseline 12.7~ns (leakage onset)\\
time step & 0.25~ns, $N_t\ge40$ & piecewise-constant propagation\\
basis modes & $\min(40,\lfloor N_t/3\rfloor)$ per quadrature & 17 modes per quadrature at 12.7~ns ($N_t=51$)\\
amplitude bound $A$ & $2.5\,\pi/T_g$ & $2.5\times$ peak of a Hann pulse of equal area\\
filter & Butterworth, order 4, 300~MHz & zero-phase (forward--backward)\\
leakage weight $\lambda$ & 7 & in $L=1-\bar F+\lambda P_{\mathrm{leak}}$\\
detuning noise $\sigma_\delta/2\pi$ & 2.5~MHz (0.5--25~MHz in sweeps) & quasi-static\\
amplitude noise $\sigma_\gamma$ & 2\% (1--12.7\% in Stage 1) & quasi-static\\
correlation $\rho$ & 0 ($-0.99$ to $0.99$ in Stage 1) & between $\delta$ and $\gamma$\\
\bottomrule
\end{tabular}
\caption{Model and simulation parameters.}
\label{tab:params}
\end{table}

\paragraph*{Hamiltonian and dissipation.} The Hamiltonian is Eq.~(\ref{M-eq:H}),
with parameters in Table~\ref{tab:params}. With $N=4$ the model contains
the two leakage levels a fast pulse can reach; the two-level model is the
same code with the matrices truncated, so that leakage is impossible by
construction. Dissipation enters through
\begin{equation}
\dot\varrho=-i[H,\varrho]+\sum_k\Big(C_k\varrho C_k^\dagger-\tfrac12\{C_k^\dagger C_k,\varrho\}\Big),
\qquad C_1=\sqrt{\gamma_1}\,a,\;\; C_2=\sqrt{2\gamma_\varphi}\,a^\dagger a,
\label{eq:lindblad}
\end{equation}
with $\gamma_1=1/T_1$ and $\gamma_\varphi=1/T_2^*-1/(2T_1)$. A 12.7~ns gate
loses $10^{-4}$--$10^{-3}$ of fidelity to these processes, almost
independently of the control and identically for every quasi-static draw.
On the qubit, $\delta$ adds $\tfrac\delta2Z$ (a rotation about $z$) and
$\gamma$ scales the drive (a rotation about the drive axis), so the two noise
sources rotate the qubit about different axes.

\paragraph*{Pulse parameterization.} Each quadrature is a sum of Hann-windowed
sine modes chosen so that the pulse starts and ends at zero:
\begin{equation}
\tilde u_I(t)=\sum_{n=1}^{M}c^{I}_n\,\sin\!\frac{(2n-1)\pi t}{T_g}\,h(t),\qquad
\tilde u_Q(t)=\sum_{n=1}^{M}c^{Q}_n\,\sin\!\frac{2n\pi t}{T_g}\,h(t),
\label{eq:basis}
\end{equation}
with the Hann window $h(t)=\tfrac12(1-\cos 2\pi t/T_g)$. The in-phase modes
are symmetric about $T_g/2$ and the quadrature modes antisymmetric, the
natural symmetry of a DRAG-type correction \cite{Motzoi2009}, although the
optimizer is free to combine them arbitrarily. The raw envelope is bounded
smoothly,
\begin{equation}
u=A\tanh(\tilde u/A),
\label{eq:tanh}
\end{equation}
which keeps the pulse inside the amplitude budget $\pm A$ with an exact chain
rule. An earlier version clipped the pulse and back-propagated as if it had
not (pitfall~4 in Sec.~\ref{sec:pitfalls}). Because
$A=2.5\,\pi/T_g$ scales with the gate time, the budget is equally tight at
every $T_g$ and the only physical change along a gate-time sweep is the gate
time itself. The optimized controls use the whole budget: at 12.7~ns the raw
envelopes of both quadratures saturate at $\pm A/2\pi=\pm98$~MHz, and the
filtered pulses that reach the qubit peak at about $115$~MHz because of the
filter's overshoot (Fig.~\ref{fig:pulses}).

\paragraph*{Bandwidth filter.} Real arbitrary-waveform generators and cabling
limit the bandwidth. The bounded envelope is passed through a 4th-order
Butterworth low-pass filter with a 300~MHz cutoff, applied forward and
backward (zero phase, Gustafsson initial conditions). Because the filter is
linear in the samples, we build it once as an exact $N_t\times N_t$ matrix $F$;
the simulated pulse is $u_f=Fu$ and the gradient with respect to $u$ is
$F^{\mathsf T}\,\partial L/\partial u_f$.

\paragraph*{Loss.} For each noise draw the loss is Eq.~the loss of Sec.~\ref{M-sec:model},
evaluated on four input states. Each input is a mixture of the ideal pure state
and the thermal population of the leakage levels,
$\varrho_{\mathrm{in}}=(1-p_{\mathrm{bg}})|\psi\rangle\langle\psi|+P\varrho_{\mathrm{th}}P$,
with $P$ the projector onto levels $\ge2$ and $p_{\mathrm{bg}}=6.8\times10^{-5}$,
so that a perfect gate still shows the thermal leakage background. The
per-state figure of merit is
$J_\psi=\mathrm{Re}\,\mathrm{Tr}\big[(\varrho_{\mathrm{tgt}}-\lambda P)\,\varrho_{\mathrm{out}}\big]$
with $\varrho_{\mathrm{tgt}}=U_{\mathrm{target}}\varrho_{\mathrm{in}}U_{\mathrm{target}}^\dagger$,
and $L=1-\tfrac14\sum_\psi J_\psi$. This is linear in the output state, which
makes the gradient an adjoint (backward) propagation.

\paragraph*{Noise ensembles.} Each draw is a pair $(\delta,\gamma)$. Three
disjoint ensembles share one distribution: training ($K=128$, seed 11),
validation (300, seed 22) and test (1000, seed 33). A replication seed $s$
shifts all three seeds by $1000s$. Correlated channels are generated from
independent normals $\delta$ and $g$ as
\begin{equation}
\gamma=\rho\,\sigma_\gamma\,\frac{\delta}{\sigma_\delta}+\sqrt{1-\rho^2}\,g ,
\qquad g\sim\mathcal N(0,\sigma_\gamma^2),
\label{eq:corr}
\end{equation}
which gives $\mathrm{corr}(\delta,\gamma)=\rho$ with unchanged marginals and
reproduces the uncorrelated draws exactly at $\rho=0$, so every Stage-0 point
remains valid. Using one distribution for all three ensembles matters: an
earlier version trained on wider noise than it tested on, which by itself
made CVaR controls look worse (pitfall~2).


\subsection{Two simulators}
\label{sec:simulators}
\paragraph*{Full Lindblad model.} The reference simulator propagates the
density matrix under Eq.~\eqref{eq:lindblad} with piecewise-constant
controls. On each time step the Liouvillian $\mathcal L_k$ (a
$N^2\times N^2$ matrix, $16\times16$ for 4 levels) is exponentiated exactly.
The gradient with respect to the control on step $k$ needs the derivative of
$e^{\mathcal L_k\Delta t}$ in the direction $\partial\mathcal L_k/\partial u$,
which we obtain exactly from the block-matrix identity \cite{VanLoan1978}
\begin{equation}
\exp\!\begin{pmatrix}\mathcal L_k & \partial\mathcal L_k\\ 0 & \mathcal L_k\end{pmatrix}\!\Delta t
=\begin{pmatrix}e^{\mathcal L_k\Delta t} & \int_0^{\Delta t}e^{\mathcal L_k(\Delta t-s)}\,\partial\mathcal L_k\,e^{\mathcal L_ks}\,ds\\ 0 & e^{\mathcal L_k\Delta t}\end{pmatrix}.
\label{eq:vanloan}
\end{equation}
Combined with a forward pass for the state and a backward pass for the
adjoint of $J$, this gives the exact gradient at the cost of about three
simulations. Per-draw evaluations are distributed over worker processes, and
draw $i$ is always returned at row $i$, which keeps the bootstrap paired
(pitfall~3).

\paragraph*{Fast closed-system model.} Optimizing with the full model costs
minutes per run. $T_1$ and $T_2^*$ act nearly identically on every
quasi-static draw, so to a good approximation they add a constant to each
per-draw loss, and $\cvar_\alpha(L+c)=\cvar_\alpha(L)+c$ does not move the
optimum. We therefore \emph{optimize} on a closed-system model in which the
Hamiltonians of all $K$ draws and all $N_t$ time steps are diagonalized in one
vectorized call, $H_{k,j}=V\,\mathrm{diag}(e)\,V^\dagger$, so each step
propagator is $V\,\mathrm{diag}(e^{-ie\Delta t})V^\dagger$. The derivative of
the propagator in direction $X$ follows from the Daleckii--Krein formula
\cite{Higham2008}: in the eigenbasis,
\begin{equation}
\big[D\,e^{-iH\Delta t}[X]\big]_{mn}=G_{mn}\,(V^\dagger XV)_{mn},\qquad
G_{mn}=\begin{cases}\dfrac{e^{-ie_m\Delta t}-e^{-ie_n\Delta t}}{e_m-e_n}, & e_m\ne e_n,\\[6pt]
-i\Delta t\,e^{-ie_m\Delta t}, & e_m=e_n.\end{cases}
\label{eq:dk}
\end{equation}
With dissipation switched off, the fast model reproduces the Lindblad code to
$10^{-7}$ in both loss and gradient (self-test). It is fast enough to run a full multistart point in about a minute.

\paragraph*{Every reported number is a full-Lindblad re-evaluation.} After
optimization, all saved controls (mean, CVaR at each $\alpha$, reverse starts)
are re-simulated in the full Lindblad model on the test ensemble, and all
statistics are recomputed there. How well the two simulators agree is
quantified in Sec.~\ref{sec:simcheck}.


\subsection{Objectives and optimizer}
All objectives share one code path. For a control with parameters
$\mathbf c$, the training losses $L_k(\mathbf c)$ and their gradients are
computed for all $K$ draws and combined with weights $w_k$:
\begin{equation}
\mathcal J(\mathbf c)=\sum_{k=1}^{K}w_k\,L_k(\mathbf c),\qquad
\nabla\mathcal J=\sum_{k=1}^{K}w_k\,\nabla L_k(\mathbf c).
\label{eq:objective}
\end{equation}
The mean uses $w_k=1/K$. CVaR uses the exact Rockafellar--Uryasev weights of
Eq.~(\ref{M-eq:ru}): sort the draws by loss, give the worst
$\lfloor\alpha K\rfloor$ weight $1/(\alpha K)$ and the next one the fractional
remainder. At $\alpha=0.05$ and $K=128$ that is six draws at weight $0.156$ and
one at $0.0625$. Where two draws tie at the boundary the weights are a valid
subgradient, and ties do not occur in practice with continuous noise.

\paragraph*{Optimizer.} L-BFGS-B \cite{Byrd1995} with at most 500 iterations
(1000 in the strong search), $\mathrm{ftol}=10^{-12}$ and
$\mathrm{gtol}=10^{-8}$, identical for every objective. The CVaR objective is
only piecewise smooth, because the set of tail draws changes as the control
moves; in practice L-BFGS-B handles this well, and the in-sample audit
(below) checks that it actually decreases the training CVaR.

\paragraph*{Restarts.} Each objective gets the same number of restarts, 4 by
default (20 in the strong search). The list is: the Hann pulse of the correct
area; random perturbations of it (independent Gaussian noise of scale $0.3A$
on every sample, projected back onto the basis); and, along a sweep, a
continuation from the previous point's mean control, rescaled to keep the
rotation area. For CVaR, one restart is replaced by the warm start from
$\umean$.


\subsection{Extracting the error rotation vector}

Given the $2\times2$ qubit block $U_q$ of a propagator, form
$E=U_{\mathrm{target}}^\dagger U_q$. If $E$ is not unitary (leakage), replace
it by its polar factor $E\to U\!V^\dagger$ from the singular value
decomposition $E=U\Sigma V^\dagger$, the closest unitary in Frobenius norm.
Divide by $\sqrt{\det E}$ to land in $SU(2)$, where
$E=c\,I-i\,\mathbf v\cdot\boldsymbol\sigma$ with
$c=\tfrac12\mathrm{Re}\,\mathrm{Tr}\,E$ and
$v_j=-\tfrac12\mathrm{Im}\,\mathrm{Tr}(E\sigma_j)$. The square root leaves a
sign ambiguity $E\leftrightarrow-E$, resolved by choosing $c\ge0$, i.e.\ a
rotation angle $\epsilon\le\pi$. Then $\epsilon=2\,\mathrm{atan2}(|\mathbf v|,c)$
and $\boldsymbol\epsilon=\epsilon\,\mathbf v/|\mathbf v|$.

\subsection{Exact gradient of the closed-system loss}

For one draw and one input, $J=\mathrm{Re}\,\mathrm{Tr}[O\,\varrho_N]$ with
$\varrho_{k+1}=U_k\varrho_kU_k^\dagger$. Define the backward-propagated
observable $\Lambda_{N}=O$, $\Lambda_k=U_k^\dagger\Lambda_{k+1}U_k$. A change
$dU_k$ of one step propagator changes $J$ by
$dJ=2\,\mathrm{Re}\,\mathrm{Tr}\big[\Lambda_{k+1}\,dU_k\,\varrho_kU_k^\dagger\big]$
(for Hermitian $O$). With $M_k=\varrho_kU_k^\dagger\Lambda_{k+1}$ and
$dU_k=V\big(G\circ(V^\dagger XV)\big)V^\dagger\,du$ from Eq.~(\ref{eq:dk}),
\[
\frac{\partial J}{\partial u_k}=2\,\mathrm{Re}\sum_{m,n}G_{mn}\,(V^\dagger XV)_{mn}\,(V^\dagger M_kV)_{nm},
\]
with $X=(1+\gamma)(a+a^\dagger)/2$ for $u_I$ and $(1+\gamma)\,i(a^\dagger-a)/2$
for $u_Q$. One forward pass stores the $\varrho_k$, one backward pass builds
the $\Lambda_k$, and the cost is independent of the number of control
parameters. The gradient with respect to the basis coefficients follows by
the chain rule through the filter $F$ and the bound (\ref{eq:tanh}),
$\partial u/\partial\tilde u=1-(u/A)^2$.

\subsection{The four-state loss metric}
\label{sec:W}

Let the error $E=U_{\mathrm{target}}^\dagger U=\cos(\epsilon/2)I-i\sin(\epsilon/2)\,\mathbf n\cdot\boldsymbol\sigma$
act on an input with Bloch vector $\mathbf r$. The state fidelity is
$F=|\langle\psi|E|\psi\rangle|^2=|\cos(\epsilon/2)-i\sin(\epsilon/2)\,\mathbf n\cdot\mathbf r|^2
=\cos^2(\epsilon/2)+\sin^2(\epsilon/2)(\mathbf n\cdot\mathbf r)^2$, so
\[
1-F_{\mathbf r}=\sin^2(\epsilon/2)\big(1-(\mathbf n\cdot\mathbf r)^2\big).
\]
The four inputs $|0\rangle,|1\rangle,|{+}\rangle,|{+}i\rangle$ have
$\mathbf r=\hat z,-\hat z,\hat x,\hat y$, and
$\sum_{\mathbf r}(\mathbf n\cdot\mathbf r)^2=2n_z^2+n_x^2+n_y^2=1+n_z^2$. Hence
\[
1-\bar F=\tfrac14\sin^2(\epsilon/2)\,(3-n_z^2)
\approx\tfrac1{16}\big(3\epsilon_x^2+3\epsilon_y^2+2\epsilon_z^2\big)
=\boldsymbol\epsilon^{\mathsf T}W\boldsymbol\epsilon,\qquad W=\tfrac1{16}\mathrm{diag}(3,3,2).
\]
We checked this expression against the four-state loss computed by the
closed-system simulator for the 2-level model. On the transmon the thermal
admixture of the inputs and the leakage term add small corrections that $W$
does not include.


% =====================================================================
\section{Fairness protocol, statistics and pipeline errors}
\label{sec:fairnessall}
\subsection{The fairness protocol}
\label{sec:fairness}

A small difference between two optimized controls is easily created by
the optimization procedure rather than by the objectives. The protocol below
was built up in response to specific failures (Sec.~\ref{sec:pitfalls}).
\begin{enumerate}[leftmargin=*,itemsep=2pt]
\item \emph{Equal restarts and tolerances} for every objective.
\item \emph{Warm start.} One CVaR restart begins at the selected mean control.
L-BFGS-B only accepts steps that decrease the objective, so this run
\emph{cannot} end with a higher training CVaR than $\umean$. If it does, the
code is wrong. This \emph{in-sample audit} is checked at every point.
\item \emph{Budget parity.} The warm-started CVaR run has received the mean
optimization plus one more run. The mean objective therefore gets one extra
continuation run from its own optimum, kept if it improves the validation
mean. This removes the possibility that CVaR ``wins'' only because it was
optimized for longer.
\item \emph{Selection on validation data only.} Each objective picks its
winning restart by its own risk measure on the validation ensemble. The test
ensemble is used only for the reported numbers.
\item \emph{Reverse warm start.} For each CVaR control, the mean optimizer is
restarted from $\ucvar$. The result is classified, at a relative tolerance
of $10^{-3}$ in validation mean loss, as: \emph{same basin} (it walks back to
the mean loss of $\umean$, so CVaR shifted the solution within one basin);
\emph{different basins} (it stays at a worse mean loss); or \emph{artifact}
(it finds a better mean control than the multistart did, so the multistart
had missed a basin).
\item \emph{Conservative comparator.} The reported $\Rt$ compares $\ucvar$
with the best mean-objective control found at that point by any route
(multistart, continuation or reverse start), chosen on validation data. The
mean side thus receives more search than the CVaR side. A positive $\Rt$ is a
\emph{lower bound} on the CVaR advantage. The converse does not hold: if a
reverse start upgraded the comparator by more than $0.1\%$, a negative $\Rt$
may just mean that the CVaR search did not get the same upgrade. Such cells
are drawn with open markers, their same-budget value is listed in
Sec.~\ref{sec:fulltable}, and no cell is dropped.
\end{enumerate}


\subsection{Evaluation and statistics}
\label{sec:stats}

The outcome metrics, on the test ensemble in the full Lindblad model, are
the tail gap and the relative tail gain,
\begin{equation}
D(\alpha)=\cvar_\alpha(L[\umean])-\cvar_\alpha(L[\ucvar]),\qquad
\Rt(\alpha)=\frac{D(\alpha)}{\cvar_\alpha(L[\umean])}.
\label{eq:rtail}
\end{equation}
$\Rt>0$ means the CVaR control has the better tail; $\Rt=0.2$ means its
worst-5\% average error is 20\% lower. Uncertainty comes from a \emph{paired}
bootstrap \cite{Efron1993}: the 1000 test draws are resampled with
replacement 2000 times, the \emph{same} resampled indices are used for both
controls, and $D$ and $\Rt$ are recomputed each time; the 95\% interval is
the 2.5--97.5 percentile range. Pairing matters because the two controls are
evaluated on the same draws, and much of the draw-to-draw variation is
common to both; an unpaired bootstrap would give much wider intervals. A
point counts as a significant divergence when the interval excludes zero, and
as within noise when it contains zero. The \emph{out-of-sample audit} for the
warm-started CVaR control requires $D\ge0$ up to its interval.

\paragraph*{Test-set resolution.} With 1000 test draws, $\cvar_{0.05}$ averages
50 draws, which resolves differences of a few percent. $\cvar_{0.01}$
averages 10 test draws but was optimized on $1.28$ training draws, which is why
that level overfits (Sec.~\ref{sec:audit}). Replication across independent
noise seeds shows that out-of-sample $\Rt$ scatters by about $\pm0.04$ at
points where the objectives are equivalent, which sets the practical
resolution of the study.


\subsection{Audit and replication}
\label{sec:audit}
Before any physics, the pipeline must pass its own test: a CVaR optimizer
started from $\umean$ must not end with a worse training CVaR, and should not
end with a significantly worse test CVaR either.

\begin{figure}[htbp]
\centering
\figorbox{figures/fig_audit.pdf}{0.5\textwidth}
\caption{Audit (E1, E6): $\Rt$ at the leakage onset, $T_g=12.7$~ns, for five
CVaR levels and three independent noise seeds (``fast'' is the first seed,
from the main run), full Lindblad, CVaR control against the mean control it
was warm-started from. The scatter between seeds ($\approx\pm0.04$) sets the
resolution of the study. At $\alpha=0.01$ the CVaR rests on $1.3$ training
draws and overfits.}
\label{fig:audit}
\end{figure}

\begin{table}[htbp]
\centering\small
\begin{tabular}{@{}lccccc@{}}
\toprule
$\alpha$ & 0.01 & 0.025 & 0.05 & 0.10 & 0.20\\
\midrule
seed 0: $\Rt$ & $-0.038$ & $-0.000$ & $-0.010$ & $+0.006$ & $-0.004$\\
\quad 95\% CI & $[-0.147,0.062]$ & $[-0.012,0.012]$ & $[-0.031,0.008]$ & $[-0.005,0.018]$ & $[-0.016,0.009]$\\
seed 1: $\Rt$ & $-0.137$ & $-0.041$ & $-0.034$ & $-0.038$ & $-0.010$\\
seed 2: $\Rt$ & $+0.016$ & $+0.013$ & $+0.018$ & $+0.005$ & $-0.001$\\
\bottomrule
\end{tabular}
\caption{E1 and E6, $T_g=12.7$~ns, full Lindblad, test ensemble. The
in-sample audit passes for every seed and $\alpha$. As the audit requires,
$\Rt$ here compares the CVaR control with the mean control it was warm-started
from (equal budget). Out of sample, $\Rt$
scatters across seeds by about $\pm0.04$ around zero; with seed 1 the CVaR
control is significantly \emph{worse} at every $\alpha$.}
\label{tab:audit}
\end{table}

All five CVaR levels pass the in-sample audit, and none fails out of sample
(Table~\ref{tab:audit}, Fig.~\ref{fig:audit}). At $\alpha=0.01$ the CVaR averages just
$0.01\times128\approx1.3$ training draws; its training CVaR improves while
its test CVaR does not, which is finite-sample overfitting rather than a
physical effect. We therefore treat $\alpha=0.01$ as below the resolution of
$K=128$.

The replication (E6) sharpens this. With three independent training, validation
and test ensembles, $\Rt$ at the onset ranges from $-0.04$ to $+0.02$ for
$\alpha\ge0.025$. There is no systematic CVaR advantage, and one seed shows a
significant disadvantage. This is expected: CVaR at $\alpha=0.05$ is estimated
from about 6 of the 128 training draws, so the CVaR control fits its training
sample more closely than the mean control does. Where the theorem says the two
objectives share a minimizer, CVaR can only add estimation noise.


\subsection{Seven ways to manufacture a spurious result}
\label{sec:pitfalls}

Each of the following occurred in an earlier version of this pipeline and was
caught by the audit. They are listed because each would have produced a
publishable-looking but false conclusion, and because most of them are
generic to comparisons between optimization objectives.

\begin{table}[htbp]
\centering\small
\begin{tabular}{@{}p{0.03\textwidth}p{0.30\textwidth}p{0.29\textwidth}p{0.28\textwidth}@{}}
\toprule
& Error & Spurious effect & Cure\\
\midrule
1 & loss defined as $-(F-\lambda P_{\mathrm{leak}})<0$ & sign of $\Rt$ flipped; a worse control read as better & loss $\ge0$, larger is worse\\
2 & training noise wider than test noise & CVaR pays for robustness never tested & one distribution for all ensembles\\
3 & parallel results collected in completion order & ``paired'' bootstrap silently unpaired & draw $i$ always at row $i$\\
4 & hard clipping with unclipped gradient & wrong gradients on up to 58\% of samples & smooth bound $A\tanh(\tilde u/A)$\\
5 & auxiliary regularizer 1--5$\times$ the loss & $\Rt\approx0.3$ at the onset, $\approx0$ without & regularizer removed\\
6 & mean optimizer stuck in a poor basin & any other objective looks good & reverse warm start, best-mean comparator\\
7 & mean given extra search, CVaR not & CVaR looks arbitrarily bad ($\Rt=-3.6$) & conservative-comparator rule\\
\bottomrule
\end{tabular}
\caption{Pipeline errors found and removed.}
\label{tab:pitfalls}
\end{table}

\begin{enumerate}[leftmargin=*,itemsep=1pt]
\item Defining the loss as $-(F-\lambda P_{\mathrm{leak}})<0$ flips the sign
of $\Rt$, because its denominator becomes negative: a CVaR control that was
\emph{worse} on every test draw read as ``better''.
\item Training on wider noise (2.5~MHz, 2\%) than the test set (1~MHz, 1\%)
makes a CVaR control pay, in nominal fidelity, for robustness against tails
that never occur in the test set.
\item Collecting parallel results in completion order silently unpairs a
``paired'' bootstrap: the same index then refers to different draws for the
two controls.
\item Hard-clipping the pulse while back-propagating as if unclipped gives
wrong gradients on up to 58\% of samples, since that fraction of samples sat
on the rail.
\item An auxiliary filter-function regularizer, sized as 2\% of the fidelity
\emph{reward}, was $1$--$5\times$ the infidelity \emph{loss}. The optimizer
then mostly minimized a noise model that was neither in the ensemble nor in
the reported CVaR. It manufactured $\Rt\approx0.3$ at the leakage onset;
without it, $\Rt\approx0$.
\item A mean optimizer stuck in a poor basin makes \emph{any} other objective
look good. The reverse warm start exposed 2--9\% better mean controls at
long gate times.
\item The converse. A mean control found with \emph{extra} search, if CVaR is
not given the same, can make CVaR look arbitrarily bad: in one cell
($T_g=20$~ns, $\alpha=0.2$), $\Rt=-3.6$ against a mean control found by a
reverse start, whereas the same cell reads $+0.11$ at equal budget. Hence the
conservative-comparator rule of Sec.~\ref{sec:fairnessall}: positive values are
lower bounds, and negative values in upgraded cells are not interpreted.
\end{enumerate}
Each has the same cure: an audit that could have made the headline result
fail. Items 6 and 7 are two sides of one problem, the asymmetry of search
effort between objectives, and they show that neither the equal-budget nor
the best-found comparison is safe alone. Reporting both, as
Sec.~\ref{sec:fulltable} does, is the honest option.


% =====================================================================
\section{Error-geometry diagnostics}
\label{sec:diagnostics}

To understand \emph{why} two controls differ, we extract the full error
of every test draw, not only its loss.

\paragraph*{Error rotation vectors.} For each draw, the qubit block of the
closed-system propagator is compared with the target,
$E=U_{\mathrm{target}}^\dagger U_{2\times2}$. If leakage makes the block
slightly non-unitary, it is replaced by its nearest unitary (the polar factor
from a singular value decomposition). After removing the global phase,
$E=\cos(\epsilon/2)I-i\sin(\epsilon/2)\,\mathbf n\cdot\boldsymbol\sigma$, and
the error rotation vector $\boldsymbol\epsilon=\epsilon\mathbf n$ follows from
the trace and the Pauli components of $E$, with the sign fixed so that
$\epsilon\le\pi$. For a 2-level model this is exact; its four-state loss is
$\boldsymbol\epsilon^{\mathsf T}W\boldsymbol\epsilon$ to leading order, which
we verified against the simulator (Sec.~\ref{sec:W}).

\paragraph*{Whitened noise coordinates.} The two noise variables are
standardized, $x_1=\delta/\sigma_\delta$, $x_2=\gamma/\sigma_\gamma$, and
rotated into independent standard normals:
\begin{equation}
z_\pm=\frac{x_1\pm x_2}{\sqrt{2(1\pm\rho)}} ,
\label{eq:whiten}
\end{equation}
ordered so that the direction with the larger variance $1+|\rho|$ comes first
($z_+$ for $\rho\ge0$, $z_-$ for $\rho<0$). With one active channel,
$z=x_1$ or $x_2$. Whitening keeps the fit well conditioned at $|\rho|=0.99$,
where $x_1$ and $x_2$ are almost collinear.

\paragraph*{Expansion and $\kappa$.} Over the 1000 test draws we fit, by least
squares, each component of $\boldsymbol\epsilon$ as a quadratic polynomial in
$\mathbf z$:
\begin{equation}
\boldsymbol\epsilon(\mathbf z)\approx\mathbf e_0+\sum_iA_{\cdot i}z_i+\sum_{i\le j}Q_{\cdot ij}z_iz_j .
\label{eq:fit}
\end{equation}
The coherent error is $\mathbf e_0$, the first-order response $A$ (a
$3\times n$ matrix) and the second-order coefficients $Q$ (a $3\times
n(n+1)/2$ array). The mechanism variable is
$\kappa=\|Q\|_F/\|A\|_F$, computed for the mean control. For one channel
it reduces to $|\mathbf q|/|\mathbf a|$, where $\boldsymbol\epsilon\approx\mathbf e_0+\mathbf az+\mathbf qz^2$.

\paragraph*{Reduced-model prediction.} From the fitted $A$ and $Q$ of the mean
control, $R^{\mathrm{model}}_{0.05}$ (Proposition~\ref{M-prop:reduced}) is
computed with no free parameter: we draw $6\times10^4$ standard normal vectors
$\mathbf z$, form the noise part $A\mathbf z+Q(\mathbf z)$, and minimize
$\cvar_{0.05}$ of $L=(\mathbf e_0+A\mathbf z+Q(\mathbf z))^{\mathsf T}W(\cdots)$
over $\mathbf e_0\in\mathbb R^3$ by Nelder--Mead, starting from the exact mean
optimum $\mathbf e_0=-\E\,Q(\mathbf z)$. Then
$R^{\mathrm{model}}=1-\cvar^{\min}/\cvar(\mathbf e_0^{\mathrm{mean}})$.

\paragraph*{Offsets.} For points where one noise direction dominates, the
coherent offset of each control is summarized as
$c=-\mathbf e_0\cdot\mathbf q/|\mathbf q|^2$, the offset in units of the
second-order error, to be compared with Remark~\ref{M-rem:precomp}.

\paragraph*{Loss-quadratic geometry.} For Proposition~\ref{M-prop:twoaxis}, the
per-draw loss is also fitted directly as
$L\approx L_0+\mathbf g^{\mathsf T}\mathbf x+\mathbf x^{\mathsf T}M\mathbf x$;
the weights $\lambda_1\ge\lambda_2$ are the eigenvalues of
$C^{1/2}MC^{1/2}$ with $C$ the channel correlation matrix, and the loss
anisotropy is $(\lambda_1-\lambda_2)/(\lambda_1+\lambda_2)$. This fit is
coarser than Eq.~\eqref{eq:fit}, because the loss is quartic in the noise once
$Q$ matters: its $R^2$ ranges from $0.55$ to $0.998$ over the diagnosed
mean controls.


% =====================================================================
\section{Simulator validation}
\label{sec:simcheck}
\begin{figure}[htbp]
\centering
\figorbox{figures/fig_simcheck.pdf}{0.95\textwidth}
\caption{How the fast optimizer model relates to the full Lindblad model, at
all 44 points with saved error clouds. (a) The absolute tail gap $D$ agrees
across the two models over four decades (median ratio $1.000$, interquartile
range $0.998$--$1.022$). (b) The ratio $\Rt$ does not, at low losses: the
Lindblad value (vertical) is predicted to within $0.005$ by the fast $D$
divided by the fast CVaR plus the $T_1/T_2$ floor (orange; median error
$0.0003$), while the fast $\Rt$ alone (crosses) overstates it where the
losses fall below the floor.}
\label{fig:simcheck}
\end{figure}

The fast model leaves out $T_1$ and $T_2^*$. To check that this does not bias
the conclusions, we compared both models at every point for which
closed-system error clouds were saved (44 points, Fig.~\ref{fig:simcheck}).
Two results emerge. First, the \emph{absolute} gap $D$ is the same in both
models: the median ratio $D_{\mathrm{Lindblad}}/D_{\mathrm{fast}}$ is $1.000$,
with interquartile range $0.998$--$1.022$. Dissipation adds almost the same
constant to both controls' losses and cancels in the difference, as
translation invariance of CVaR predicts. Second, the \emph{relative} gain
$\Rt$ is diluted, because the constant enters its denominator:
\begin{equation}
\Rt^{\mathrm{Lindblad}}\approx\frac{D_{\mathrm{fast}}}{\cvar_{\alpha}^{\mathrm{fast}}(L[\umean])+\ell_{\mathrm{floor}}} ,
\label{eq:dilution}
\end{equation}
where $\ell_{\mathrm{floor}}$, the mean-loss difference between the two
models for $\umean$, ranges from $1.4\times10^{-4}$ to $1.06\times10^{-3}$.
Equation~\eqref{eq:dilution} predicts the Lindblad $\Rt$ with median absolute
error $0.0003$ and maximum $0.005$. The fast $\Rt$ alone has median error
$0.003$ but reaches $0.19$ at points whose losses lie well below the floor.
All conclusions in this paper are drawn from the Lindblad values, which are
the conservative ones; where the dilution matters (the prospective test at a
quarter of the noise) both are reported.


% =====================================================================
\section{Experiments}
\label{sec:experiments}

\begin{table}[htbp]
\centering\small
\begin{tabular}{@{}llp{0.50\textwidth}r@{}}
\toprule
ID & Name & What is varied & Points\\
\midrule
E1 & Audit & $\alpha\in\{0.01,0.025,0.05,0.1,0.2\}$ at $T_g=12.7$~ns, 4 levels & 1\\
E2 & Leakage sweep & $T_g=6.3$--$32.5$~ns across the leakage onset; 4 and 2 levels & 9+5\\
E3 & Channel split & detuning only vs amplitude only, $T_g=12.7$--$20.8$~ns & 3+3\\
E4 & Phase-variance sweep & $\sigma_\delta/2\pi=0.5$--$25$~MHz at $T_g=8$, $12.7$, $20$~ns; 2 levels, detuning only & 16\\
E5 & Strong search & 20 restarts, 1000 iterations, $T_g=15.4$, $20.8$, $32.5$~ns; 4 and 2 levels & 3+3\\
E6 & Replication & E1 and two E4 points with two further noise seeds & 6\\
E7 & Several channels & $\sigma_\gamma=0.01$--$0.127$ and $\rho=-0.99$ to $0.99$ at $T_g=12.7$~ns, 2 levels; replication seeds & 16\\
E8 & Prospective test & four new points at $\rho=0.99$, predictions fixed in advance & 4\\
E9 & Dominance test & candidate pools of 376--491 controls; E4's phase variances at 8, 12.7, 20~ns (2 levels, 3 seeds at 12.7~ns) and 12.7, 20.8~ns (4 levels) & 42\\
\bottomrule
\end{tabular}
\caption{Experiments. Every point comprises one mean optimization and CVaR
optimizations at $\alpha=0.05$ and $0.2$ (all five levels in E1), each with
its own restarts and reverse start, all re-evaluated in the full Lindblad
model.}
\label{tab:experiments}
\end{table}

The experiments are listed in Table~\ref{tab:experiments}. The
\emph{leakage onset} used to place the E2 sweep is the gate time at which the
nominal leakage of a Hann pulse exceeds $5\times$ the thermal background,
found by log--log interpolation on 30 gate times between 3 and 120~ns
(Fig.~\ref{fig:leak}). E4 matches $\sigma_\phi^2=(\sigma_\delta T_g)^2$ across
the three gate times, so that the \emph{collapse test} can check whether
$\sigma_\phi$ alone controls the result. E7 runs at a detuning noise that
gives $\sigma_\phi^2=0.04$, inside the single-channel equivalence region, so
that any divergence must come from the second channel or from the
correlation. At $\sigma_\gamma=0.127$ the amplitude channel has the same phase
variance as the detuning channel. E8 is described in Sec.~\ref{M-sec:mechanism} and E9 in Sec.~\ref{M-sec:dominance}.

\paragraph*{Computation.} All simulations were run on a personal laptop
(10 CPU cores, 8 worker processes, Python with NumPy and SciPy). A
two-level Stage-1 point (one mean and two CVaR multistarts with reverse
starts) takes 70--90~s with the fast model, and a full-Lindblad re-evaluation of
one point about 14~s.

% =====================================================================
\section{Additional results}
\label{sec:addresults}
\subsection{Leakage onset}
\label{sec:leakonset}
\begin{figure}[htbp]
\centering
\figorbox{figures/leak_vs_Tg.pdf}{0.48\textwidth}
\caption{Nominal leakage of a Hann pulse versus gate time. The onset
($5\times$ the thermal background of $6.8\times10^{-5}$) is at
$T_g^{\mathrm{on}}=12.7$~ns, and at $11.6$ or $13.4$~ns for $10\times$ or
$3\times$, so the onset does not hinge on the threshold chosen.}
\label{fig:leak}
\end{figure}

Leakage of a Hann pulse falls from 12\% at 3~ns to the thermal floor by
$\sim$20~ns (Fig.~\ref{fig:leak}). The onset, $T_g^{\mathrm{on}}=12.7$~ns,
defines the leakage-dominated region $T_g<T_g^{\mathrm{on}}$.


\subsection{Gate-time sweep and strong search}
\label{sec:sweep}
\begin{table}[htbp]
\centering\footnotesize\setlength{\tabcolsep}{3.2pt}
\begin{tabular}{@{}lccccccccc@{}}
\toprule
$T_g$ (ns) & 6.3 & 7.4 & 8.5 & 9.9 & 11.5 & 12.7 & 15.4 & 20.8 & 32.5\\
Hann leakage & $3\!\times\!10^{-2}$ & $2\!\times\!10^{-2}$ & $8\!\times\!10^{-3}$ & $3\!\times\!10^{-3}$ & $1\!\times\!10^{-3}$ & $5\!\times\!10^{-4}$ & $2\!\times\!10^{-4}$ & $7\!\times\!10^{-5}$ & $7\!\times\!10^{-5}$\\
\midrule
$\Rt$, 4 levels & $-0.003$ & $-0.011$ & $-0.016$ & $-0.017$ & $-0.021$ & $-0.022$ & $+0.104$ & $+0.102$ & $+0.064$\\
$\Rt$, 2 levels & $-0.026$ & -- & -- & -- & -- & $-0.026$ & $-0.012$ & $-0.007^{\dagger}$ & $+0.009$\\
\midrule
strong search, 4 lev. & & & & & & & $+0.100$ & $+0.117^{*}$ & $+0.059$\\
strong search, 2 lev. & & & & & & & $-0.008^{*}$ & $-0.021$ & $-0.004$\\
\bottomrule
\end{tabular}
\caption{E2 and E5, $\alpha=0.05$, full Lindblad, CVaR control against the
best mean control. Below and at the onset, $|\Rt|\le0.026$ in every case.
Strong search: 20 restarts per objective and 1000 iterations. $^{*}$A
reverse start improved the mean control by more than 0.1\%, and the value
shown is against that improved control (two-level 15.4~ns: $+0.076$ at equal
budget). 95\% CIs of the strong search:
$[0.077,0.122]$, $[0.092,0.138]$, $[0.036,0.080]$. $^{\dagger}$Against the
multistart mean control alone this cell reads $+0.098$; a reverse start found
a mean control that matches the CVaR control's tail, so the apparent
divergence was a missed basin (pitfall~6, Sec.~\ref{sec:pitfalls}).}
\label{tab:sweep}
\end{table}


Throughout the leakage-dominated region, $6.3\le T_g\le12.7$~ns, the mean and
CVaR controls coincide within $|\Rt|\le0.026$, with 4 levels and with 2
(Table~\ref{tab:sweep}). Where $\Rt$ differs from zero at all it is slightly
negative, the same estimation-variance penalty seen in the audit. The CVaR optimizer started from $\umean$ moves
the pulse by only 1--2\% and gains under 1\% in-sample. This is not because
leakage is absent from the problem: a Hann pulse at 6.3~ns leaks 3\%. Both
optimizers suppress leakage to near the thermal floor
($\approx7$--$9\times10^{-5}$), and the residual loss distributions remain
ordered, exactly as Corollary~\ref{M-cor:general} requires.


\subsection{Channel split}
\label{sec:channelsplit}
\begin{table}[htbp]
\centering\small
\begin{tabular}{@{}lccc@{}}
\toprule
$T_g$ (ns) & 12.7 & 15.4 & 20.8\\
\midrule
both channels & $-0.022$ & $+0.104$ & $+0.102$\\
amplitude only & $-0.013$ & $-0.166^{\ddagger}$ & $-0.001$\\
detuning only & $-0.015$ & $+0.074$ & $+0.177$\\
\bottomrule
\end{tabular}
\caption{E3, 4 levels, $\alpha=0.05$, full Lindblad, against the best mean
control. Amplitude noise alone never produces a positive divergence.
$^{\ddagger}$Comparator asymmetry: the best mean control here came from a
reverse start off the $\alpha=0.2$ CVaR control, a basin the $\alpha=0.05$
CVaR search never visited; at equal budget the cell reads $-0.007$.}
\label{tab:channel}
\end{table}

Amplitude noise alone gives no CVaR advantage at any gate time
(Table~\ref{tab:channel}); its one large negative cell is a comparator
artefact, not a physical effect. This rules out an alternative mechanism in which
the level-2 AC Stark shift converts amplitude error into a quadratic phase.
Any divergence requires detuning noise.

The strong search (E5) settles whether the long-gate divergence is real.
With 20 restarts and 1000 iterations per objective, and CVaR compared against
the best mean control found anywhere, the divergence persists: $\Rt=0.100$, $0.117$ and $0.059$ at
15.4, 20.8 and 32.5~ns (Table~\ref{tab:sweep}). At 15.4 and 32.5~ns the
reverse start returns to $\umean$, so the CVaR solution is a genuine shift
within the same basin.

\paragraph{The transmon diverges earlier than the qubit.} A two-level qubit
with exactly the same noise shows no divergence at any gate time up to
32.5~ns. With the same 20-restart, 1000-iteration search as the transmon,
$\Rt=-0.008$, $-0.021$ and $-0.004$ at 15.4, 20.8 and 32.5~ns
($\sigma_\phi^2=0.06$, $0.11$, $0.26$), where the transmon gives $0.10$,
$0.12$ and $0.06$. The 4-restart runs agree ($-0.012$, $-0.007$, $+0.009$).
At 15.4~ns the qubit shows once more how an apparent divergence arises: at
equal budget $\Rt=+0.076$, but a reverse start finds a mean control with the
same tail. In absolute terms, at 20.8~ns the transmon's
tail gap is $D=1.2\times10^{-3}$, against $-4\times10^{-5}$ for the qubit. The
higher transmon levels therefore lower the divergence threshold and amplify
it once detuning noise is present, while by themselves (amplitude noise only)
they produce none.


\subsection{Phase-variance sweep and replication}
\label{sec:sigmasweep}
\begin{table}[htbp]
\centering\small
\begin{tabular}{@{}lcccccc@{}}
\toprule
$\sigma_\delta/2\pi$ (MHz) & 0.5 & 1 & 2.5 & 5 & 8 & 16\\
$\sigma_\phi^2$ & 0.0016 & 0.0064 & 0.040 & 0.16 & 0.41 & 1.6\\
\midrule
$\Rt(0.05)$ & $-0.001^{\dagger}$ & $-0.002$ & $+0.002$ & $+0.029$ & $+0.077$ & $+0.275$\\
95\% CI & & $[-0.005,0.001]$ & $[-0.003,0.007]$ & $[0.026,0.032]$ & $[0.068,0.085]$ & $[0.247,0.298]$\\
$\Rt(0.20)$ & $-0.026^{\dagger}$ & $+0.000$ & $+0.002$ & $+0.013$ & $+0.034$ & $+0.135$\\
\bottomrule
\end{tabular}
\caption{E4, 2 levels, detuning noise only, $T_g=12.7$~ns, full Lindblad.
The reverse warm start returns to $\umean$ at 5, 8 and 16~MHz (same basin),
so these are genuine in-basin shifts. $^{\dagger}$Losses $\sim10^{-5}$,
at the optimizer's precision floor; the comparator was upgraded by a reverse
start.}
\label{tab:sigma}
\end{table}



On a bare qubit with detuning noise alone, the setting closest to
Theorem~\ref{M-thm:main}, $\Rt$ is zero within its confidence interval for
$\sigma_\phi^2\le0.04$. This is direct numerical confirmation of the theorem
in its stated regime (Sec.~\ref{sec:proofs}). Beyond
$\sigma_\phi^2\approx0.1$ it grows smoothly and monotonically, reaching
$0.275$ at $\sigma_\phi^2=1.6$ (Table~\ref{tab:sigma}). At every divergent
point the mean optimizer, restarted from $\ucvar$, walks back to $\umean$:
the two optima sit in one basin, and CVaR genuinely shifts the solution
within it.

\paragraph{Replication.} Repeating two points with two new noise seeds
(E6): at $\sigma_\phi^2=1.6$, $\Rt(0.05)=0.275$, $0.271$ and $0.314$ for the
three seeds, all with confidence intervals far from zero. At
$\sigma_\phi^2=0.16$, $\Rt(0.05)=0.029$ and $0.036$ for two seeds. For the
third, the multistart comparison gives $+0.050$, but a reverse start from the
$\alpha=0.2$ CVaR control found a mean control with a better $\alpha=0.05$ tail
than the $\alpha=0.05$ CVaR search reached ($\Rt=-0.40$ against it): a CVaR
search failure of the kind described in pitfall~7, not a reversal of the
effect.


\subsection{Several channels}
\label{sec:stage1}
\begin{table}[htbp]
\centering\footnotesize\setlength{\tabcolsep}{3.5pt}
\begin{tabular}{@{}lcccccc@{}}
\toprule
\multicolumn{7}{@{}l}{\emph{Amplitude width, $\rho=0$}}\\
$\sigma_\gamma$ & 0.01 & 0.02 & 0.04 & 0.08 & 0.127 & 0.127, no detuning\\
$\Rt(0.05)$ & $-0.134^{\circ}$ & $-0.029$ & $-0.013$ & $+0.044$ & $+0.017$ & $+0.084^{\circ}$\\
95\% CI & $[-0.16,-0.11]$ & $[-0.06,0.00]$ & $[-0.03,0.01]$ & $[0.01,0.08]$ & $[-0.01,0.04]$ & $[-0.04,0.18]$\\
$\kappa$ & 0.08 & 0.10 & 0.12 & 0.23 & 0.53 & --\\
\midrule
\multicolumn{7}{@{}l}{\emph{Correlation, $\sigma_\gamma=0.127$}}\\
$\rho$ & $-0.99$ & $-0.9$ & $-0.5$ & $+0.5$ & $+0.9$ & $+0.99$\\
$\Rt(0.05)$ & $+0.226$ & $+0.167$ & $+0.114$ & $+0.013$ & $+0.077$ & $+0.289$\\
95\% CI & $[0.13,0.30]$ & $[0.08,0.24]$ & $[0.05,0.17]$ & $[-0.01,0.04]$ & $[0.07,0.09]$ & $[0.24,0.34]$\\
$\kappa$ & 1.98 & 1.12 & 0.66 & 0.47 & 0.83 & 1.73\\
\bottomrule
\end{tabular}
\caption{E7, 2 levels, $T_g=12.7$~ns, full Lindblad, against the best mean
control. $\kappa=|Q|/|A|$ for the mean control. Replication seeds: at
$\sigma_\gamma=0.08$, $+0.028$ $[0.01,0.05]$ and $+0.003$ $[-0.02,0.02]$; at
$\rho=+0.99$, $+0.073$ $[-0.08,0.20]$ and $+0.334$ $[0.29,0.37]$.
$^{\circ}$Comparator upgraded by a reverse start.}
\label{tab:stage1}
\end{table}


\paragraph{P1 and P2: the two-axis mechanism is real but weak.} With both
channels independent, the objectives stay within a few percent of each other
over the whole amplitude range (Table~\ref{tab:stage1}). The one significant
positive point, $\sigma_\gamma=0.08$, replicates weakly ($+0.044$, $+0.028$,
$+0.003$ over three seeds). The geometric signature Proposition~\ref{M-prop:twoaxis}
predicts, a more isotropic CVaR control, appears in the first seed (loss
anisotropy $A=0.87\to0.79$) but not in the other two ($0.75\to0.74$,
$0.88\to0.89$). The two-axis mechanism is therefore allowed by the theory but
contributes little in practice: optimized controls can change the split of
their error between the axes only slightly. Amplitude noise alone at
$\sigma_\gamma=0.127$ gives no significant divergence either.

\begin{figure}[htbp]
\centering
\figorbox{figures/fig_clouds.pdf}{\textwidth}
\caption{Error clouds: the error rotation vector of each of the 1000 test
draws, projected onto its two principal axes and coloured by loss, for the
mean and $\cvar_{0.05}$ controls at $\sigma_\gamma=0.127$. Cross: zero
error. Circle: the average error of the cloud. Left pair ($\rho=0$): the
cloud is two-dimensional, curved by second-order error, and the two controls
are almost identical. Right pair ($\rho=0.99$): the linear spread is
cancelled along one direction and the cloud collapses to a thin,
one-sided streak. The CVaR control shifts the whole streak (circle to the
left of the cross), so that its long right-hand tail is pulled back towards
zero at the cost of the typical draw.}
\label{fig:clouds}
\end{figure}

\paragraph{P3's condition fails, and the reason is informative.} Read
unconditionally, as it was first framed, P3 fails: correlating the channels
does not restore equivalence but produces the largest divergences of the
whole study, $\Rt=0.23$ at $\rho=-0.99$ and $0.29$ at $\rho=+0.99$, with every
reverse start returning to $\umean$. The error clouds show why. With
correlated noise the pulse need only cancel one combined first-order error
direction instead of two, and it does: the first-order response along the
combined noise direction $x_1\pm x_2$ (the one that dominates as
$|\rho|\to1$) falls from $0.076$ at $\rho=0$ to $0.013$--$0.017$ at
$|\rho|=0.99$, and the mean loss of $\umean$ drops fourfold. What remains is
second order in the single noise variable, so the linear-coupling condition
of P3 fails, with $\kappa$ above 1. Figure~\ref{fig:clouds} shows the
difference directly: the independent-channel cloud is a broad two-dimensional
fan, whereas the correlated-channel cloud is a narrow, one-sided streak whose
extent is set by the second-order term. Rank-one \emph{noise} is not rank-one
\emph{linear error}: the optimization itself removes the linearity that
equivalence needs. One of three seeds at $\rho=+0.99$ is not significant
($+0.073$, with a wide interval); the other two are.

\begin{figure}[htbp]
\centering
\figorbox{figures/fig_pulses.pdf}{0.9\textwidth}
\caption{Mean and $\cvar_{0.05}$ controls at $\sigma_\gamma=0.127$ (filtered
envelopes as applied to the qubit). Top two rows: the in-phase and quadrature
envelopes, visually indistinguishable at this scale. Bottom row: their
difference. With independent channels (left) the CVaR control differs by
0.8\% in norm; with correlated channels (right) by 6.6\%, with an oscillating
correction of a few MHz, the change that produces a 29\% lower tail
infidelity.}
\label{fig:pulses}
\end{figure}

\paragraph{How different are the pulses?} Divergent CVaR controls are small
modifications of the mean controls (Fig.~\ref{fig:pulses}). At the strongest
Stage-1 divergence the relative change in the control envelopes is 6.6\%,
against 0.8\% at a non-divergent point; in the gate-time sweep (E2) the CVaR optimizer moved
the pulse by 1--2\% in the equivalence region. Wherever the reverse start
returns to $\umean$, the two objectives select nearby points in one basin,
not different pulse families. The effect of the
small change on the error is nonetheless large and systematic, as the next
section shows.


\label{sec:pulses}

\subsection{Independence of the experimental parameters}
\label{sec:independence}
\begin{table}[htbp]
\centering\small
\begin{tabular}{@{}lccc@{}}
\toprule
covariate & marginal & given $\kappa$ & given $R^{\mathrm{model}}$\\
\midrule
$|\rho|$ & $+0.70$ & $-0.00$ ($p=0.98$) & $+0.08$ ($p=0.69$)\\
detuning phase variance $\sigma_\phi^2$ & $+0.09$ & $+0.24$ ($p=0.22$) & $-0.03$ ($p=0.90$)\\
amplitude phase variance & $+0.53$ & $-0.19$ ($p=0.32$) & $+0.28$ ($p=0.15$)\\
total phase variance & $+0.41$ & $+0.25$ ($p=0.20$) & $+0.10$ ($p=0.60$)\\
gate time $T_g$ & $-0.29$ & $-0.38$ ($p=0.05$) & $-0.15$ ($p=0.45$)\\
\bottomrule
\end{tabular}
\caption{Does $\kappa$ carry the information of the experimental knobs?
Spearman correlation of $\Rt(0.05)$ with each covariate, marginal and partial
(controlling for $\kappa$ or $R^{\mathrm{model}}$), 28 two-level points,
including the prospective test. Conversely, $\kappa$ keeps a partial
correlation of $0.61$--$0.84$ ($p\le6\times10^{-4}$) and $R^{\mathrm{model}}$
of $0.78$--$0.89$ ($p\le10^{-6}$) against every covariate.}
\label{tab:independence}
\end{table}

\paragraph{Independence of the knobs.} Correlation $|\rho|$ predicts $\Rt$ on
its own (Spearman $0.70$), but once $\kappa$ is known it adds nothing (partial
$-0.00$, $p=0.98$); the same holds for noise strength
(Table~\ref{tab:independence}). Gate time keeps a borderline residual given
$\kappa$ ($-0.38$, $p=0.05$, one of ten tests), driven by the 8~ns point, which
diverges strongly at a moderate $\kappa=0.82$. It disappears given
$R^{\mathrm{model}}$ ($p=0.45$): the extra information is in the error
geometry, which $\kappa$ discards and the reduced model keeps. Within these data, $\rho$, $T_g$ and noise strength
affect the divergence only through $\kappa$ and the error geometry, as
Proposition~\ref{M-prop:reduced} requires. This test is observational: the
covariates were not varied independently of one another, and the two-level
data cover $T_g$ only between 6.3 and 32.5~ns.


\subsection{Pre-compensation offsets}

The offsets of Fig.~\ref{M-fig:precomp}(c), per point (coherent offset
$c=-\mathbf e_0\cdot\mathbf q/|\mathbf q|^2$ of each control, in units of its
second-order error): detuning scan at 5, 8 and 16~MHz: mean control
$0.51$, $0.51$, $0.49$; $\cvar_{0.2}$ $1.61$, $1.62$, $1.71$;
$\cvar_{0.05}$ $2.46$, $2.48$, $2.54$. Correlated channels, $\cvar_{0.05}$:
$2.00$ ($\rho=-0.99$), $1.87$ ($-0.9$), $1.02$ ($+0.9$), $1.98$ and $2.52$
($+0.99$, two seeds), and $1.63$, $1.94$, $2.03$, $2.46$ at the four points of
the prospective test; $\cvar_{0.2}$ $1.15$--$1.80$; mean controls
$0.73$--$0.94$. The exception $\rho=+0.9$ retains a large second noise
direction.

% =====================================================================
\section{Dominance test: details}
\label{sec:dominancedetails}
\paragraph*{Pools and statistics.} The candidate pool of a setting contains a
geometrically thinned subset of the L-BFGS-B iterates of every run (mean,
CVaR at $\alpha=0.05$ and $0.2$, reverse starts) plus every run's final
control, deduplicated. Every candidate is evaluated with the closed-system
simulator on a fresh ensemble of 4000 draws (seed $44+1000s$), split into two
halves of 2000. Statistics are computed over the competitive candidates (mean
loss within a factor of two of the pool's best). The full-Lindblad
$\Rt$ of the standard comparison is computed as in the rest of the study.
A confirmed reversal requires a separation of more than 2\% in both mean and
$\cvar_{0.05}$ on both halves, with the same orientation on both halves. The
pool regret selects the best-mean and best-CVaR candidates on one half,
scores both by $\cvar_{0.05}$ on the other,
$R_{\mathrm{pool}}=1-\cvar(\text{best-CVaR})/\cvar(\text{best-mean})$, and
averages the two assignments.

\paragraph*{Transmon and second basin.} With four levels the same pattern
holds, shifted. Reversals stay near 1\% at the smallest noise instead of zero,
and the shape ratio is 1.8--3.1 instead of 5.58, because a noise-independent
leakage floor makes up 56--84\% of the mean loss there. At 20.8~ns the
transmon crosses over earlier than the qubit (4.6\% reversals at
$\sigma_\phi^2=0.04$, 22\% at 0.16). At 8 and 20~ns and $\sigma_\phi^2=1.6$
the competitive pool is dominated by a second basin, found by the reverse
starts, with a four to five times lower mean loss but a heavier relative tail
($\cvar_{0.05}/\text{mean}\approx10$).

\paragraph*{The unrestricted crossing test.}
\label{sec:crossing}
Table~\ref{tab:crossing} lists, for every two-level setting with
$\sigma_\phi^2\le0.04$, the crossing statistic as first specified (any
crossing of the quantile functions by more than 2\%, no mean separation
required), together with the location of the crossings and the fraction of
crossing pairs in which the control with the lower mean has the larger
noise-free loss $L_0$. The crossings are all located in the lower half of the
distribution, and in all but one seed they are driven by the coherent error.
The lowest quantiles of a $\chi^2_1$ loss are tiny ($0.4\%$ of the mean at
$q=0.05$), so a residual coherent error of a few tenths of a percent of the
mean decides them, whereas the upper half, which determines both the mean
and the CVaR, is set by the noise. The restricted statistic of the main text
(pairs with at least 10\% mean separation, crossings at $q\ge0.5$) is zero at
all of these settings.

\begin{table}[htbp]
\centering\small
\begin{tabular}{@{}rrrrrrrr@{}}
\toprule
$T_g$ (ns) & seed & $\sigma_\phi^2$ & crossing pairs & fraction & median $q$ & max $q$ & lower mean has larger $L_0$\\
\midrule
8 & 0 & 0.0016 & 250 & 0.113 & 0.06 & 0.19 & 0.96\\
8 & 0 & 0.0064 & 988 & 0.233 & 0.07 & 0.25 & 0.91\\
8 & 0 & 0.04 & 1343 & 0.268 & 0.19 & 0.45 & 0.91\\
12.7 & 0 & 0.0016 & 257 & 0.133 & 0.07 & 0.42 & 0.98\\
12.7 & 1 & 0.0016 & 141 & 0.192 & 0.10 & 0.36 & 0.96\\
12.7 & 2 & 0.0016 & 45 & 0.045 & 0.05 & 0.09 & 0.49\\
12.7 & 0 & 0.0064 & 194 & 0.119 & 0.08 & 0.27 & 0.84\\
12.7 & 1 & 0.0064 & 277 & 0.210 & 0.07 & 0.34 & 0.91\\
12.7 & 2 & 0.0064 & 107 & 0.103 & 0.07 & 0.12 & 0.29\\
12.7 & 0 & 0.04 & 663 & 0.215 & 0.17 & 0.41 & 0.96\\
12.7 & 1 & 0.04 & 1232 & 0.308 & 0.15 & 0.46 & 0.93\\
12.7 & 2 & 0.04 & 677 & 0.240 & 0.12 & 0.43 & 0.78\\
20 & 0 & 0.0016 & 64 & 0.062 & 0.08 & 0.16 & 0.97\\
20 & 0 & 0.0064 & 151 & 0.146 & 0.08 & 0.27 & 0.95\\
20 & 0 & 0.04 & 469 & 0.270 & 0.15 & 0.44 & 0.93\\
\bottomrule
\end{tabular}
\caption{Unrestricted crossing test at small noise, two levels. ``median
$q$'' and ``max $q$'': median and maximum over crossing pairs of the quantile
level of the largest crossing.}
\label{tab:crossing}
\end{table}

\begin{table}[htbp]
\centering\small
\begin{tabular}{@{}lrrr@{}}
\toprule
family & $\tau$ & reversals & $R_{\mathrm{pool}}$\\
\midrule
ordered, $\sin^2(\sigma_iZ/2)$ & $1$ & $0$ & $0$\\
pre-compensation, $(Z^2-c_i)^2$ & $-0.46$ & $0.76$ & $0.42$\\
\bottomrule
\end{tabular}
\caption{Synthetic validation of the statistics (\texttt{stage2.py selftest}), run before the production run.}
\label{tab:synthetic}
\end{table}

% =====================================================================
\section{The controlled experiment, leakage-aware $\kappa$, and generality}
\label{sec:stage3}

\subsection{Design and calibration}
\label{sec:stage3design}

$\kappa$ is a property of the mean-optimal control and is therefore known only
after an optimization, so a design that holds it fixed has to solve for the
noise level that produces it. For each cell $(\kappa^\star,\rho,T_g)$ of the
grid $\kappa^\star\in\{0.20,0.35,0.55\}$, $\rho\in\{0,0.5,0.9\}$,
$T_g\in\{8,12.7,20\}$~ns we scale the base widths
$(\sigma_\delta,\sigma_\gamma)=(2.5~\mathrm{MHz},0.127)$ by one number $s$ and
solve $\kappa(s)=\kappa^\star$ by secant iteration in $\log s$, seeded with the
empirical scaling $\kappa\propto s^{1/2}$ and stopped at
$|\log\kappa/\kappa^\star|<\log 1.05$ or six evaluations. Each evaluation is a
complete mean multistart plus a fit of Eq.~(\ref{M-eq:expansion}) on the test
ensemble. The 27 cells took 112 mean optimizations; 21 converged inside the
5\% tolerance, and the worst cell ended 25\% away.

Each calibrated cell is then run with the standard protocol of
Sec.~\ref{sec:fairnessall} at two noise seeds, and the calibrated control is offered
to the mean side as one of its restarts (the CVaR side is warm-started from
the selected mean control as always, so this keeps the comparison
conservative). The run's mean control is nevertheless better optimized than
the calibration's and cancels more linear error, so its $\kappa$ is larger:
by a median factor $1.08$ over the 54 points, $1.30$ at $T_g=8$~ns, and up to
$1.74$ in one cell. The achieved ranges are $0.18$--$0.33$, $0.35$--$0.56$ and
$0.51$--$0.74$ for the three rows, which overlap slightly; all analyses use the
measured $\kappa$ of the control that was run.

\paragraph*{Statistics.} Because a cell can return $\Rt\le0$, ratios to a
prediction are not usable (a clipped ratio turns "no gain" into an arbitrarily
large log-residual), and with three cells per knob value a Kruskal--Wallis test
has little power. The tests reported in the main text are therefore partial
rank correlations of $\Rt$ with a knob, controlling for $\kappa$ or for $\Rm$,
computed over all 54 points and over the 36 whose mean loss exceeds the
$T_1/T_2$ floor of $3\times10^{-4}$. Below that floor the CVaR estimate is
dominated by its six training draws: six of the eighteen $\rho=0.9$ cells
return a significantly negative $\Rt$.

\subsection{Leakage-aware error coordinates}
\label{sec:leakmetric}

Let $U$ be the full $N$-level propagator of a draw and
$E=U_{\mathrm{target}}^\dagger U$ with the target embedded as the identity on
levels $\ge2$. Write $E_q$ for the qubit block and $l_{kj}=E_{kj}=U_{kj}$,
$k\ge2$, $j\in\{0,1\}$ for the leakage amplitudes. Taking the polar factor
$W_u$ of $E_q$, normalizing it to $SU(2)$ and applying the \emph{same} global
phase to $l$ (a gauge choice, since the loss is invariant under it) gives the
coordinate
$x=(\bm\epsilon;\,\mathrm{Re}\,l;\,\mathrm{Im}\,l)\in\mathbb R^{3+4(N-2)}$.

Since $U$ is unitary, $E_q^\dagger E_q+l^\dagger l=I$, so
$E_q=W_u(I-l^\dagger l)^{1/2}\approx W_u(I-\tfrac12l^\dagger l)$ and, for an
input $\psi$ in the computational subspace,
\[
F_\psi=|\langle\psi|E_q|\psi\rangle|^2
\approx1-\tfrac14\big[\epsilon^2-(\bm\epsilon\cdot\mathbf r)^2\big]-\langle\psi|l^\dagger l|\psi\rangle .
\]
Averaging over the four inputs gives the metric of Sec.~\ref{sec:W} for the
rotation part and $\bar P_{\rm leak}=\tfrac14\sum_\psi\langle\psi|l^\dagger l|\psi\rangle$
for the leakage part; the loss $L=1-\bar F+\lambda\bar P_{\rm leak}$ therefore
carries $(1+\lambda)$ on the leakage term. With
$S=\sum_\psi|\psi\rangle\langle\psi|=2I+\tfrac12(X+Y)$ and $G=\bar S$,
\[
\bar P_{\rm leak}=\tfrac14\sum_{k\ge2}l_k^\dagger G\,l_k ,
\]
which in the real coordinates $(\mathrm{Re}\,l,\mathrm{Im}\,l)$ is the
positive-definite form quoted in the main text. \texttt{stage3.py selftest}
checks it against the exact four-state loss of random unitaries: the relative
error is $2\times10^{-5}$ at an error size of $3\times10^{-4}$ and falls as the
cube of that size, as a second-order metric must.

The mechanism variable in these coordinates is
$\kappa_{\rm lk}=\|W^{1/2}Q\|_F/\|W^{1/2}A\|_F$, and the reduced model is
evaluated in the same space, with the offset free in all coordinates. For
$N=2$ the leakage block is empty and $\kappa_{\rm lk}$ reduces to the
metric-weighted $\kappa_W$, which differs from $\kappa$ only by the weighting
$W_\epsilon=\mathrm{diag}(3,3,2)/16$.

\subsection{Generality grid}
\label{sec:stage3generality}

The target rotation is installed by rebuilding the closed-system model with
$U_{\mathrm{target}}=\exp(-i\theta\,\hat n\!\cdot\!\bm\sigma/2)$ embedded in
the $N$-level space; the amplitude budget scales with $\theta$, so
$X_\pi$ is as tightly constrained as $X_{\pi/2}$, and the Hann seed is scaled
to the correct area. Bases are (i) the $M$ Hann-windowed sine modes per
quadrature of Sec.~\ref{sec:modeldetails}, (ii) $M=6$ of them, and (iii) $M$
windowed piecewise-constant segments, which matches the parameter count of
(i) at $M=17$. Anharmonicity is varied on the four-level model only.

\begin{figure}[htbp]
\centering
\includegraphics[width=0.98\textwidth]{fig_stage3_generality.pdf}
\caption{The generality grid. (a) $\Rt$ against $\kappa$ for the three
rotation angles (two-level, 17 sine modes) and for the four-level
anharmonicity scan. (b) The pulse-basis comparison at $X_{\pi/2}$: 17 sine
modes against 17 piecewise-constant segments (the same number of parameters)
and 6 sine modes. The two piecewise-constant points are the outliers, one of
them at $\Rt=-0.31$. (c) Measured against the reduced model over all 22
points.}
\label{fig:generality}
\end{figure}

Two caveats limit what these points can be compared with. First, the Lindblad
worker has the $X_{\pi/2}$ target at $-200$~MHz compiled into it, so points
with another target or anharmonicity are closed-system only; their JSON
records this and the validation step skips them. Second, and consequently, the
best-mean comparator of Sec.~\ref{sec:fairnessall}, which the Lindblad validation
applies, had to be applied separately to the saved closed-system losses
(\texttt{stage3\_recomp.py}). It changed 46 of the 76 Stage-3 points, by more
than $0.05$ at four of them; at the two-level $X_{\pi/2}$ point with 17 modes
it removed a spurious $\Rt=+0.31$, bringing it to $-0.02$ in agreement with
the phase-variance sweep at the same setting. The numbers reported in the main
text and in Table~\ref{tab:stage3C} are the corrected ones.

\begin{table}[htbp]
\centering\scriptsize\setlength{\tabcolsep}{3.5pt}
\caption{The controlled experiment, all 54 points. $\kappa_{\rm cal}$: the
calibration control; $\kappa_{\rm run}$: the control actually run; $C$: the
coefficient of Eq.~(\ref{M-eq:kappa2}); $\Rm$: reduced model; $\Rt$:
closed-system, best-mean comparator, $\alpha=0.05$. $\sigma_\delta$ in MHz.}
\label{tab:stage3A}
\begin{tabular}{@{}cccccccccccc@{}}
\toprule
$\kappa^\star$ & $\rho$ & $T_g$ & seed & $\sigma_\delta$ & $\sigma_\gamma$ &
$\kappa_{\rm cal}$ & $\kappa_{\rm run}$ & $C$ & $\Rm$ & $\Rt$ & $\E L$\\
\midrule
0.20 & 0.0 & 8 & 0 & 2.41 & 0.122 & 0.257 & 0.310 & 1.98 & 0.123 & $+0.094$ & 2e-03 \\
0.20 & 0.0 & 8 & 1 & 2.41 & 0.122 & 0.257 & 0.280 & 3.57 & 0.164 & $+0.138$ & 2e-03 \\
0.20 & 0.0 & 12.7 & 0 & 1.26 & 0.064 & 0.206 & 0.245 & 0.22 & 0.018 & $-0.025$ & 4e-04 \\
0.20 & 0.0 & 12.7 & 1 & 1.26 & 0.064 & 0.206 & 0.236 & 0.13 & 0.011 & $-0.055$ & 4e-04 \\
0.20 & 0.0 & 20 & 0 & 1.28 & 0.065 & 0.193 & 0.194 & 0.74 & 0.023 & $-0.005$ & 8e-04 \\
0.20 & 0.0 & 20 & 1 & 1.28 & 0.065 & 0.193 & 0.178 & 1.25 & 0.032 & $+0.012$ & 8e-04 \\
0.20 & 0.5 & 8 & 0 & 1.54 & 0.078 & 0.209 & 0.274 & 1.32 & 0.087 & $-0.005$ & 5e-04 \\
0.20 & 0.5 & 8 & 1 & 1.54 & 0.078 & 0.209 & 0.271 & 1.49 & 0.091 & $+0.050$ & 6e-04 \\
0.20 & 0.5 & 12.7 & 0 & 0.83 & 0.042 & 0.207 & 0.228 & 0.07 & 0.007 & $-0.050$ & 1e-04 \\
0.20 & 0.5 & 12.7 & 1 & 0.83 & 0.042 & 0.207 & 0.224 & 0.02 & 0.004 & $-0.016$ & 1e-04 \\
0.20 & 0.5 & 20 & 0 & 0.91 & 0.046 & 0.191 & 0.206 & 0.12 & 0.006 & $-0.003$ & 2e-04 \\
0.20 & 0.5 & 20 & 1 & 0.91 & 0.046 & 0.191 & 0.216 & 0.06 & 0.003 & $+0.025$ & 2e-04 \\
0.20 & 0.9 & 8 & 0 & 1.04 & 0.053 & 0.188 & 0.327 & 0.24 & 0.071 & $-0.050$ & 8e-05 \\
0.20 & 0.9 & 8 & 1 & 1.04 & 0.053 & 0.188 & 0.327 & 0.25 & 0.072 & $-0.021$ & 8e-05 \\
0.20 & 0.9 & 12.7 & 0 & 0.41 & 0.021 & 0.208 & 0.217 & 0.01 & 0.005 & $-0.030$ & 7e-06 \\
0.20 & 0.9 & 12.7 & 1 & 0.41 & 0.021 & 0.208 & 0.205 & 0.00 & 0.003 & $-0.160$ & 8e-06 \\
0.20 & 0.9 & 20 & 0 & 0.49 & 0.025 & 0.209 & 0.227 & 0.01 & 0.007 & $+0.001$ & 2e-05 \\
0.20 & 0.9 & 20 & 1 & 0.49 & 0.025 & 0.209 & 0.224 & 0.00 & 0.005 & $+0.001$ & 2e-05 \\
0.35 & 0.0 & 8 & 0 & 3.24 & 0.165 & 0.365 & 0.394 & 3.19 & 0.219 & $+0.176$ & 4e-03 \\
0.35 & 0.0 & 8 & 1 & 3.24 & 0.165 & 0.365 & 0.390 & 4.48 & 0.211 & $+0.187$ & 4e-03 \\
0.35 & 0.0 & 12.7 & 0 & 1.97 & 0.100 & 0.357 & 0.376 & 0.12 & 0.036 & $-0.027$ & 1e-03 \\
0.35 & 0.0 & 12.7 & 1 & 1.97 & 0.100 & 0.357 & 0.371 & 0.14 & 0.040 & $-0.100$ & 1e-03 \\
0.35 & 0.0 & 20 & 0 & 2.67 & 0.136 & 0.428 & 0.367 & 0.07 & 0.012 & $-0.008$ & 4e-03 \\
0.35 & 0.0 & 20 & 1 & 2.67 & 0.136 & 0.428 & 0.401 & 0.47 & 0.059 & $+0.036$ & 4e-03 \\
0.35 & 0.5 & 8 & 0 & 2.21 & 0.112 & 0.355 & 0.379 & 2.15 & 0.171 & $+0.147$ & 1e-03 \\
0.35 & 0.5 & 8 & 1 & 2.21 & 0.112 & 0.355 & 0.373 & 1.53 & 0.150 & $+0.082$ & 1e-03 \\
0.35 & 0.5 & 12.7 & 0 & 1.52 & 0.077 & 0.357 & 0.385 & 0.02 & 0.034 & $+0.035$ & 4e-04 \\
0.35 & 0.5 & 12.7 & 1 & 1.52 & 0.077 & 0.357 & 0.415 & 0.03 & 0.050 & $+0.053$ & 5e-04 \\
0.35 & 0.5 & 20 & 0 & 1.58 & 0.080 & 0.341 & 0.351 & 0.08 & 0.022 & $-0.016$ & 8e-04 \\
0.35 & 0.5 & 20 & 1 & 1.58 & 0.080 & 0.341 & 0.352 & 0.18 & 0.046 & $-0.026$ & 9e-04 \\
0.35 & 0.9 & 8 & 0 & 1.91 & 0.097 & 0.345 & 0.561 & 2.94 & 0.326 & $+0.226$ & 4e-04 \\
0.35 & 0.9 & 8 & 1 & 1.91 & 0.097 & 0.345 & 0.555 & 3.12 & 0.331 & $+0.039$ & 4e-04 \\
0.35 & 0.9 & 12.7 & 0 & 0.70 & 0.035 & 0.347 & 0.380 & 0.00 & 0.046 & $-0.022$ & 2e-05 \\
0.35 & 0.9 & 12.7 & 1 & 0.70 & 0.035 & 0.347 & 0.383 & 0.00 & 0.047 & $-0.049$ & 2e-05 \\
0.35 & 0.9 & 20 & 0 & 0.77 & 0.039 & 0.359 & 0.370 & 0.04 & 0.053 & $-0.082$ & 4e-05 \\
0.35 & 0.9 & 20 & 1 & 0.77 & 0.039 & 0.359 & 0.371 & 0.00 & 0.037 & $+0.011$ & 5e-05 \\
0.55 & 0.0 & 8 & 0 & 6.37 & 0.324 & 0.503 & 0.742 & 3.40 & 0.327 & $+0.190$ & 2e-02 \\
0.55 & 0.0 & 8 & 1 & 6.37 & 0.324 & 0.503 & 0.744 & 3.69 & 0.317 & $+0.183$ & 2e-02 \\
0.55 & 0.0 & 12.7 & 0 & 3.20 & 0.162 & 0.556 & 0.610 & 0.13 & 0.124 & $-0.036$ & 4e-03 \\
0.55 & 0.0 & 12.7 & 1 & 3.20 & 0.162 & 0.556 & 0.559 & 0.18 & 0.115 & $+0.101$ & 4e-03 \\
0.55 & 0.0 & 20 & 0 & 3.74 & 0.190 & 0.530 & 0.521 & 0.55 & 0.084 & $+0.030$ & 9e-03 \\
0.55 & 0.0 & 20 & 1 & 3.74 & 0.190 & 0.530 & 0.508 & 1.59 & 0.169 & $+0.120$ & 9e-03 \\
0.55 & 0.5 & 8 & 0 & 4.03 & 0.205 & 0.465 & 0.656 & 3.25 & 0.306 & $+0.202$ & 5e-03 \\
0.55 & 0.5 & 8 & 1 & 4.03 & 0.205 & 0.465 & 0.716 & 1.74 & 0.294 & $+0.068$ & 5e-03 \\
0.55 & 0.5 & 12.7 & 0 & 2.08 & 0.106 & 0.533 & 0.530 & 0.02 & 0.093 & $+0.039$ & 9e-04 \\
0.55 & 0.5 & 12.7 & 1 & 2.08 & 0.106 & 0.533 & 0.574 & 0.02 & 0.099 & $+0.023$ & 9e-04 \\
0.55 & 0.5 & 20 & 0 & 2.55 & 0.129 & 0.529 & 0.528 & 0.03 & 0.055 & $+0.007$ & 3e-03 \\
0.55 & 0.5 & 20 & 1 & 2.55 & 0.129 & 0.529 & 0.578 & 0.06 & 0.083 & $+0.074$ & 3e-03 \\
0.55 & 0.9 & 8 & 0 & 2.14 & 0.109 & 0.488 & 0.600 & 4.87 & 0.356 & $+0.264$ & 5e-04 \\
0.55 & 0.9 & 8 & 1 & 2.14 & 0.109 & 0.488 & 0.613 & 2.53 & 0.337 & $+0.086$ & 6e-04 \\
0.55 & 0.9 & 12.7 & 0 & 1.07 & 0.054 & 0.560 & 0.593 & 0.01 & 0.151 & $+0.015$ & 6e-05 \\
0.55 & 0.9 & 12.7 & 1 & 1.07 & 0.054 & 0.560 & 0.577 & 0.00 & 0.152 & $-0.132$ & 6e-05 \\
0.55 & 0.9 & 20 & 0 & 1.35 & 0.068 & 0.576 & 0.628 & 0.12 & 0.197 & $+0.041$ & 2e-04 \\
0.55 & 0.9 & 20 & 1 & 1.35 & 0.068 & 0.576 & 0.609 & 0.00 & 0.150 & $-0.011$ & 2e-04 \\
\bottomrule
\end{tabular}
\end{table}

\begin{table}[htbp]
\centering\small\setlength{\tabcolsep}{4pt}
\caption{The generality grid, all 22 points; anharmonicity in MHz, basis as
described above, noise scale relative to $(2.5~\mathrm{MHz},0)$ at
$T_g=12.7$~ns. Closed-system values with the best-mean comparator.}
\label{tab:stage3C}
\begin{tabular}{@{}cccccccccc@{}}
\toprule
levels & $\theta$ & anh. & basis & scale & $\kappa$ & $C$ & $\Rm$ & $\Rt$ & $\E L$\\
\midrule
2 & $\pi$ & -200 & 17 sine & 1 & 9.573 & 7.88 & 0.426 & $+0.344$ & 1e-07 \\
2 & $\pi$ & -200 & 17 sine & 2 & 0.798 & 8.15 & 0.432 & $+0.554$ & 1e-08 \\
2 & $\pi$ & -200 & 17 sine & 4 & 1.958 & 9.43 & 0.426 & $+0.277$ & 8e-07 \\
2 & $\pi$ & -200 & 17 sine & 8 & 2.065 & 10.17 & 0.425 & $+0.335$ & 2e-05 \\
2 & $\pi/2$ & -200 & 17 p.c. & 1 & 0.048 & 2.82 & 0.007 & $-0.305$ & 1e-04 \\
2 & $\pi/2$ & -200 & 17 p.c. & 2 & 0.095 & 5.37 & 0.048 & $+0.132$ & 3e-04 \\
2 & $\pi/2$ & -200 & 6 sine & 1 & 0.076 & 4.52 & 0.025 & $+0.003$ & 2e-04 \\
2 & $\pi/2$ & -200 & 6 sine & 2 & 0.152 & 4.53 & 0.093 & $+0.028$ & 9e-04 \\
2 & $\pi/2$ & -200 & 17 sine & 1 & 0.067 & 4.53 & 0.020 & $-0.020$ & 3e-04 \\
2 & $\pi/2$ & -200 & 17 sine & 2 & 0.136 & 4.53 & 0.076 & $+0.006$ & 1e-03 \\
2 & $\pi/2$ & -200 & 17 sine & 4 & 0.266 & 4.53 & 0.233 & $+0.117$ & 5e-03 \\
2 & $\pi/2$ & -200 & 17 sine & 8 & 0.630 & 4.53 & 0.430 & $+0.392$ & 1e-02 \\
2 & $\pi/4$ & -200 & 17 sine & 1 & 0.034 & 5.26 & 0.006 & $+0.000$ & 3e-03 \\
2 & $\pi/4$ & -200 & 17 sine & 2 & 0.063 & 5.28 & 0.021 & $+0.012$ & 1e-02 \\
2 & $\pi/4$ & -200 & 17 sine & 4 & 0.130 & 5.28 & 0.080 & $+0.033$ & 5e-02 \\
2 & $\pi/4$ & -200 & 17 sine & 8 & 0.296 & 5.31 & 0.301 & $+0.031$ & 1e-01 \\
4 & $\pi/2$ & -150 & 17 sine & 1 & 0.037 & 4.59 & 0.006 & $-0.006$ & 5e-03 \\
4 & $\pi/2$ & -200 & 17 sine & 1 & 0.035 & 4.51 & 0.005 & $-0.012$ & 3e-03 \\
4 & $\pi/2$ & -300 & 17 sine & 1 & 0.050 & 4.63 & 0.011 & $+0.013$ & 2e-03 \\
4 & $\pi/2$ & -150 & 17 sine & 2 & 0.074 & 4.60 & 0.024 & $+0.014$ & 2e-02 \\
4 & $\pi/2$ & -200 & 17 sine & 2 & 0.072 & 4.50 & 0.021 & $+0.022$ & 1e-02 \\
4 & $\pi/2$ & -300 & 17 sine & 2 & 0.099 & 4.66 & 0.043 & $+0.030$ & 5e-03 \\
\bottomrule
\end{tabular}
\end{table}

% =====================================================================
\section{Colored noise}
\label{sec:stage4}

\subsection{Ornstein--Uhlenbeck sampling and normalization}

The detuning is replaced by an Ornstein--Uhlenbeck process
$d\eta=-\eta\,dt/\tau_c+\sigma_{\mathrm{ou}}\sqrt{2/\tau_c}\,dW$, with
stationary variance $\sigma_{\mathrm{ou}}^2$ and correlation
$C(t,t')=\sigma_{\mathrm{ou}}^2e^{-|t-t'|/\tau_c}$. On the control grid of
$n$ slices of width $dt$ it is sampled exactly as a stationary AR(1) chain,
\begin{equation}
\eta_{j+1}=a\eta_j+\sigma_{\mathrm{ou}}\sqrt{1-a^2}\,\xi_j,\qquad
a=e^{-dt/\tau_c},
\end{equation}
with $\eta_0\sim\mathcal N(0,\sigma_{\mathrm{ou}}^2)$ and $\xi_j$ i.i.d.\
standard normal, so no discretization error enters the noise itself. The
accumulated phase over the gate is $\phi=dt\sum_j\eta_j$, whose variance is
exactly
\begin{equation}
\mathrm{Var}(\phi)=\sigma_{\mathrm{ou}}^2\,dt^2\sum_{j,k=0}^{n-1}a^{|j-k|}
=\sigma_{\mathrm{ou}}^2\,dt^2\Big(n+2\!\!\sum_{d=1}^{n-1}\!(n-d)a^{d}\Big).
\end{equation}
We solve this for $\sigma_{\mathrm{ou}}$ so that $\mathrm{Var}(\phi)$ equals
the target $\sigma_\phi^2$ at every $\tau_c$. The quasi-static arm is the
$\tau_c\to\infty$ limit of the same construction and reproduces the sampler of
Sec.~\ref{sec:modeldetails}. Without this normalization a short-$\tau_c$ arm would simply
be a weaker-noise arm, and the whole comparison would be confounded; the
resulting $\sigma_{\mathrm{ou}}$ grows from $2.51$~MHz quasi-static to
$18.9$~MHz at $\tau_c=0.1\,T_g$, $\sigma_\phi^2=0.41$, as it must.

\subsection{Noise coordinates}

The error-geometry fit of Sec.~\ref{sec:diagnostics} expands the error in whitened
scalar noise coordinates. A trajectory has no single scalar, so we use the
first two functionals of the path that a gate can resolve: the accumulated
phase and its first time moment about the mid-gate instant,
\begin{equation}
z_1\propto dt\sum_j\eta_j,\qquad
z_2\propto dt\sum_j(t_j-\bar t)\,\eta_j ,
\end{equation}
each centred over the ensemble, then $z_2$ orthogonalized against $z_1$ and
both standardized. Under quasi-static noise $z_2$ vanishes identically, and
the fit reduces to the scalar one. The reported $R^2_{\mathrm{fit}}$ is the
smallest coefficient of determination over the components of the error vector
when regressed on $(z_1,z_2)$ and their quadratic products; it measures
whether two coordinates suffice. They do not at short correlation time, which
is the honest limitation of this stage: $R^2_{\mathrm{fit}}$ falls to $0.010$
at $\tau_c=0.1\,T_g$. A third and higher moments would extend the reduction,
at the cost of a fit that needs many more draws to condition.

\subsection{Design and scoring}

The grid is $\tau_c/T_g\in\{0.1,0.3,1,3,10\}$ plus a quasi-static arm, crossed
with $\sigma_\phi^2\in\{0.04,0.16,0.41\}$ and seeds $\{0,1\}$: $36$ points,
all at $T_g=12.7$~ns on the two-level model, $17$ sine modes, with the
fairness protocol of Sec.~\ref{sec:fairnessall} unchanged. The best-mean comparator
is applied at save time rather than during validation, which is the fix for
the pitfall recorded in Sec.~\ref{sec:stage3generality}; it upgraded $20$ of the
$36$ points, by at most $0.064$.

Four predictions were registered before the run. Two were met as stated (T1,
memory independence; T4, the growth of the $z_2$ share), and two were
mis-specified in ways the run exposed.

\paragraph*{T2 was a special case.} The registered prediction was
$\cvar_\alpha/\E=5.58$ at every point. That is $\rho_\alpha(1)$, and holds
only when one noise direction reaches the gate. Proposition~\ref{M-prop:meff}
supersedes it; the quasi-static arm, where $m_{\mathrm{eff}}=1$ by
construction, returns $5.36$ and does satisfy the original form.

\paragraph*{T3 was unattainable.} The registered threshold was
$\mathrm{Spearman}(\kappa,\Rt)>0.6$. With two seeds per cell the replicate
standard deviation of $\Rt$ is $\sigma_{\mathrm{noise}}=0.0205$ against a
total $0.0256$, leaving $\sigma_{\mathrm{signal}}=0.0153$; classical
attenuation then caps any observable correlation at
$(1+\sigma^2_{\mathrm{noise}}/\sigma^2_{\mathrm{signal}})^{-1/2}=0.600$. The
threshold demanded a perfect underlying correlation. We report disattenuated
values instead: $\mathrm{Pearson}(\Rm,\Rt)=0.67$ observed, $1.12$
disattenuated, and $\mathrm{Spearman}(\kappa,\Rt)=0.44$ observed, $0.74$
disattenuated. The residual scatter about the fit is $0.0190$, or $0.93$ of
the replicate scatter.

Table~\ref{tab:stage4} lists every point. We record both re-specifications
here rather than in the main text so that the original registration stays
visible. Neither changes a conclusion in the
direction that favours the theory: T2's correction makes the shape prediction
harder, since it now must hold per control rather than as one number, and
T3's makes the mechanism claim weaker, since it becomes a statement about
consistency with a bound rather than about rank discrimination.

\begin{table*}[t]
\caption{All $36$ colored-noise points, ordered by decreasing correlation
time. $\tau_c/T_g=\infty$ is the quasi-static arm. $C$ and $\Rm$ come from the
two-coordinate reduction and should be read with $R^2_{\mathrm{fit}}$ in mind;
ratio is $\cvar_{0.05}/\E$; the last column is the comparator selected by the
best-mean rule of Sec.~\ref{sec:fairness}.}
\label{tab:stage4}
\begin{ruledtabular}
\begin{tabular}{ccccccccccc l}
$\tau_c/T_g$ & $\sigma_\phi^2$ & seed & $\kappa$ & $C$ & $C\kappa^2$ &
$\Rm$ & $\Rt$ & ratio & $z_2$ share & $R^2_{\mathrm{fit}}$ & comparator\\
\hline
10 & 0.04 & 0 & 0.065 & 4.92 & 0.0207 & 0.0204 & +0.0089 & 4.40 & 0.270 & 0.984 & mean\\
10 & 0.04 & 1 & 0.065 & 4.85 & 0.0204 & 0.0200 & -0.0431 & 4.48 & 0.273 & 0.982 & rev$_{0.2}$\\
10 & 0.16 & 0 & 0.130 & 4.93 & 0.0836 & 0.0738 & +0.0278 & 4.36 & 0.287 & 0.983 & mean\\
10 & 0.16 & 1 & 0.129 & 4.83 & 0.0800 & 0.0712 & -0.0122 & 4.43 & 0.286 & 0.981 & rev$_{0.05}$\\
10 & 0.41 & 0 & 0.208 & 4.83 & 0.2087 & 0.1557 & +0.0599 & 4.32 & 0.305 & 0.980 & rev$_{0.05}$\\
10 & 0.41 & 1 & 0.206 & 4.74 & 0.2001 & 0.1507 & +0.0258 & 4.36 & 0.306 & 0.977 & rev$_{0.2}$\\
3 & 0.04 & 0 & 0.053 & 1.83 & 0.0052 & 0.0053 & -0.0092 & 3.98 & 0.551 & 0.961 & mean\\
3 & 0.04 & 1 & 0.053 & 1.79 & 0.0051 & 0.0052 & +0.0020 & 4.14 & 0.543 & 0.961 & mean\\
3 & 0.16 & 0 & 0.105 & 1.69 & 0.0185 & 0.0193 & -0.0092 & 3.96 & 0.555 & 0.944 & mean\\
3 & 0.16 & 1 & 0.104 & 1.81 & 0.0196 & 0.0201 & -0.0043 & 4.12 & 0.549 & 0.961 & rev$_{0.2}$\\
3 & 0.41 & 0 & 0.165 & 1.28 & 0.0348 & 0.0413 & -0.0028 & 3.95 & 0.586 & 0.962 & rev$_{0.05}$\\
3 & 0.41 & 1 & 0.164 & 1.35 & 0.0362 & 0.0418 & -0.0031 & 4.07 & 0.578 & 0.957 & mean\\
1 & 0.04 & 0 & 0.041 & 0.33 & 0.0006 & 0.0007 & -0.0372 & 4.21 & 0.703 & 0.853 & rev$_{0.2}$\\
1 & 0.04 & 1 & 0.041 & 0.36 & 0.0006 & 0.0005 & -0.0072 & 4.42 & 0.750 & 0.908 & rev$_{0.05}$\\
1 & 0.16 & 0 & 0.077 & 0.13 & 0.0007 & 0.0013 & +0.0000 & 4.33 & 0.771 & 0.909 & mean\\
1 & 0.16 & 1 & 0.080 & 0.15 & 0.0009 & 0.0013 & -0.0070 & 4.40 & 0.761 & 0.931 & rev$_{0.05}$\\
1 & 0.41 & 0 & 0.121 & 0.04 & 0.0006 & 0.0024 & -0.0027 & 4.29 & 0.790 & 0.883 & mean\\
1 & 0.41 & 1 & 0.123 & 0.08 & 0.0012 & 0.0035 & -0.0083 & 4.34 & 0.776 & 0.890 & mean\\
0.3 & 0.04 & 0 & 0.022 & 2.63 & 0.0012 & 0.0013 & -0.0032 & 3.65 & 0.602 & 0.390 & rev$_{0.2}$\\
0.3 & 0.04 & 1 & 0.044 & 1.57 & 0.0030 & 0.0032 & -0.0173 & 3.63 & 0.581 & 0.422 & rev$_{0.2}$\\
0.3 & 0.16 & 0 & 0.044 & 1.90 & 0.0037 & 0.0037 & +0.0027 & 3.63 & 0.604 & 0.389 & mean\\
0.3 & 0.16 & 1 & 0.064 & 1.87 & 0.0077 & 0.0078 & -0.0380 & 3.61 & 0.540 & 0.495 & mean\\
0.3 & 0.41 & 0 & 0.075 & 1.47 & 0.0083 & 0.0080 & +0.0077 & 3.58 & 0.607 & 0.395 & mean\\
0.3 & 0.41 & 1 & 0.093 & 1.82 & 0.0159 & 0.0146 & -0.0323 & 3.57 & 0.543 & 0.504 & rev$_{0.2}$\\
0.1 & 0.04 & 0 & 0.018 & 3.23 & 0.0011 & 0.0013 & -0.0050 & 3.83 & 0.652 & 0.196 & rev$_{0.2}$\\
0.1 & 0.04 & 1 & 0.069 & 0.19 & 0.0009 & 0.0009 & -0.0077 & 3.44 & 0.669 & 0.264 & rev$_{0.05}$\\
0.1 & 0.16 & 0 & 0.039 & 0.76 & 0.0012 & 0.0013 & -0.0147 & 3.65 & 0.789 & 0.062 & rev$_{0.2}$\\
0.1 & 0.16 & 1 & 0.074 & 0.11 & 0.0006 & 0.0008 & -0.0017 & 3.49 & 0.797 & 0.092 & mean\\
0.1 & 0.41 & 0 & 0.057 & 0.34 & 0.0011 & 0.0013 & -0.0077 & 3.77 & 0.844 & 0.010 & rev$_{0.2}$\\
0.1 & 0.41 & 1 & 0.112 & 0.36 & 0.0046 & 0.0042 & +0.0174 & 3.33 & 0.704 & 0.213 & rev$_{0.2}$\\
$\infty$ & 0.04 & 0 & 0.074 & 4.53 & 0.0247 & 0.0240 & +0.0040 & 5.17 & 0.000 & 1.000 & mean\\
$\infty$ & 0.04 & 1 & 0.075 & 4.53 & 0.0252 & 0.0248 & -0.0236 & 5.36 & 0.000 & 1.000 & mean\\
$\infty$ & 0.16 & 0 & 0.150 & 4.53 & 0.1023 & 0.0914 & +0.0282 & 5.16 & 0.000 & 0.999 & rev$_{0.05}$\\
$\infty$ & 0.16 & 1 & 0.147 & 4.54 & 0.0985 & 0.0892 & -0.0511 & 5.37 & 0.000 & 0.999 & rev$_{0.2}$\\
$\infty$ & 0.41 & 0 & 0.244 & 4.52 & 0.2684 & 0.2049 & +0.0716 & 5.11 & 0.000 & 0.996 & mean\\
$\infty$ & 0.41 & 1 & 0.240 & 4.53 & 0.2606 & 0.1996 & +0.0386 & 5.22 & 0.000 & 0.995 & rev$_{0.05}$\\
\end{tabular}
\end{ruledtabular}
\end{table*}

% =====================================================================
\section{Channel geometry}
\label{sec:stage1b}

\subsection{Design and calibration}
The noise enters through two physically distinct channels: a static detuning
$\delta$ of standard deviation $\sigma_\delta$ and a multiplicative
drive-amplitude error of standard deviation $\sigma_a$, independent. We
parameterize them by a mixing angle, $\sigma_\delta=s\cos\theta$ and
$\sigma_a=s\sin\theta/\,\Omega_0$ with $\Omega_0$ the nominal Rabi rate, and
step $\theta$ through nine values from $0^\circ$ to $90^\circ$ in
$11.25^\circ$ increments. At each $\theta$ the overall scale $s$ is
recalibrated by a secant search so that the mean loss of the mean-optimal
control matches a common target; without that step the two channels would
differ in both geometry and strength and the comparison would be confounded.
Each probe of the search is a single short optimization ($120$ iterations)
rather than a full multistart, with one full multistart run at the calibrated
scale. Two seeds, $T_g=12.7$~ns, two-level model: $18$ points. The achieved
mean losses are $3.1$--$8.8\times10^{-4}$ for $\theta<90^\circ$; at
$\theta=90^\circ$ the control nulls the amplitude response almost entirely and
the calibration cannot hold, giving $5$--$6\times10^{-6}$.

\subsection{Pre-registered predictions and outcomes}
\begin{description}
\item[B1] The measured tail-to-mean ratio follows
$\rho_\alpha(m_{\mathrm{eff}})$ as the channel split moves $m_{\mathrm{eff}}$,
with $\mathrm{Spearman}>0.8$ and median relative error below $10\%$.
\emph{Outcome: fails as written, passes when blocked on seed.} Pooled
$\mathrm{Spearman}=+0.712$ ($p=9\times10^{-4}$), median error $7.2\%$. Within
seeds, $+0.976$ ($p=3\times10^{-5}$) and $+0.952$ ($p=3\times10^{-4}$) at
median errors $7.9\%$ and $3.4\%$; seed-centred over the $16$ in-regime
points, $+0.962$ ($p=3\times10^{-9}$). The pooled statistic is defeated by a
constant offset of $+0.563\pm0.141$ between the seeds against a prediction
that varies by only $1.51$ across the whole range of $m_{\mathrm{eff}}$. The
rank correlation of $m_{\mathrm{eff}}$ with the measured ratio is $-0.976$ and
$-0.952$ against a predicted $-1$ exactly.
\item[B2] Where the split moves $m_{\mathrm{eff}}$ by more than $0.3$, the CDFs
of the mean- and CVaR-optimal controls cross. \emph{Outcome: not scored.} The
design produced one such pair, so the test has no power. We report this rather
than reporting ``0 of 1 crossings''.
\item[B3] The $\kappa^2$ envelope continues to hold, $\max\Rt<0.10$.
\emph{Outcome: fails on $18$ points, holds on the $16$ in regime.} The two
failures are the $\theta=90^\circ$ pair, at $\kappa=6.26$ and $46.85$: a bound
of the form $C\kappa^2+O(\kappa^3)$ is not a statement about $\kappa=47$.
Excluding them, all remaining points have $\kappa\le0.116$ and
$\max\Rt=0.0275$.
\item[B4] One curve $\rho_\alpha(m_{\mathrm{eff}})$ fits both this route and
the colored-noise route of Sec.~\ref{sec:stage4}; the two residual sets are
not distinguishable. \emph{Outcome: passes.} Mann--Whitney $p=0.77$ on the
$16$ in-regime points against the $18$ Stage~4 points with
$R^2_{\mathrm{fit}}\ge0.94$ ($p=0.40$ including the $\theta=90^\circ$ pair).
Medians $-1.5\%$ and $-0.1\%$; $10$th--$90$th percentiles $[-8.0,+3.9]\%$ and
$[-7.4,+2.6]\%$. Pooled over all $34$ points,
$\mathrm{Spearman}=+0.831$ ($p=1\times10^{-9}$) at a median absolute residual
of $3.7\%$.
\end{description}

\subsection{Per-point table}
$\theta$ in degrees; $\sigma_a$ is the relative amplitude error; ratio and
prediction at $\alpha=0.05$; $\Rt$ against the best mean comparator. The last
two rows are the excluded $\theta=90^\circ$ pair.
\begin{center}\footnotesize\setlength{\tabcolsep}{4pt}
\input{figures/tab_stage1b.tex}
\end{center}

% =====================================================================
\section{Trajectory-level relaxation}
\label{sec:stage5}

\subsection{An exact one-jump decomposition}
Amplitude damping is not a Gaussian perturbation of the propagator, so the
machinery of Sec.~\ref{sec:proofs} does not apply and the loss distribution
must be constructed directly. Monte Carlo wavefunction sampling would do it,
but the quantity we need --- a tail average at $\alpha=10^{-4}$ --- is
precisely the one trajectory sampling estimates worst. We instead expand in
the number of jumps and keep every branch exactly.

Let $a$ be the lowering operator, $n=a^\dagger a$ and $\Gamma=1/T_1$. On each
control slice we use the Strang-split damped map
\begin{equation}
K_k=D^{1/2}U_kD^{1/2},\qquad D=\exp(-\tfrac{1}{2}\Gamma\,dt\,n),
\label{eq:damped}
\end{equation}
so that the no-jump branch is $K_{N}\cdots K_1$. A single jump at slice $k$
sends the state to the normalized collapsed state, after which the evolution
is again no-jump. The observable for that branch is
$\mathrm{Tr}(O\,S_k\sigma S_k^\dagger)=\mathrm{Tr}\big((S_k^\dagger OS_k)\sigma\big)$
with $S_k=K_N\cdots K_{k+1}$, so all $N$ one-jump branches are obtained from a
single backward recursion that accumulates $S_k^\dagger OS_k$. The cost is one
extra pass over the time grid, and the truncation error is $O((\Gamma T)^2)$.

Two details decide whether this agrees with the Lindblad channel, and both
were caught by validation rather than by inspection.

\emph{Bayes weighting.} The four input states of the average-gate-fidelity
estimator do not have the same probability of producing a given record: the
no-jump weight is $\mathrm{Tr}(K\rho_jK^\dagger)$, which depends on $j$. A
plain average over inputs is therefore inconsistent with the channel and
leaves an $O(\Gamma)$ discrepancy. Weighting each input by
$P(j\mid r)\propto\mathrm{Tr}(K_r\rho_jK_r^\dagger)$ removes it: the residual
disagreement with a direct Lindblad solve falls from $2.3\times10^{-5}$ to
$1.2\times10^{-6}$ at $T_1=30\,\mu$s, which is $0.46\%$ of the $T_1$
contribution itself.

\emph{The jump weight.} The natural weight $\Gamma\,dt\,\langle n\rangle$ is
inconsistent with the \emph{discretized} no-jump map at $O(\Gamma\,dt)$, which
shows up as a drifting mass defect. The consistent choice is the telescoping
trace loss of the map actually used,
$m_{jk}=\mathrm{Tr}\,\rho_j(t_k)-\mathrm{Tr}\,K_k\rho_j(t_k)K_k^\dagger$,
placed on the normalized collapsed state. With it, total probability is
conserved to $0.8\,(\Gamma T)^2$ with no floor in $dt$. Across the whole study
the largest residual mass defect is $2.2\times10^{-6}$.

\emph{Validation.} Six self-tests run before every production job: exactness
of the atomic $\cvar$ against a sorted brute-force reference including the
boundary-atom split ($1.6\times10^{-15}$); the $\Gamma$-scaling of the
discrepancy against a Lindblad reference; mass conservation against
$(\Gamma T)^2$; agreement of $\rho_\alpha$ and $\Xi_\alpha$ with their
closed forms; the mixture identity of Eq.~(\ref{M-eq:mixture}) on a two-channel
toy; and the backward recursion against a direct $O(N^2)$ construction
($6.7\times10^{-14}$). The last of these caught an off-by-one in the
recursion --- the update $S_k^\dagger OS_k$ must be applied \emph{before}
storing, since $S_k=S_{k+1}U_k$ --- which had produced a maximum relative
error of $39$.

\subsection{Design and scoring}
The grid crosses $T_g\in\{12.7,40\}$~ns, $T_1\in\{10,30\}\,\mu$s and
$\sigma_\phi^2\in\{0.04,0.16\}$ with two seeds on the two-level model, with
$\sigma_\phi^2=0.16$ run at $T_g=40$~ns only, plus a four-level arm at
$T_g=12.7$~ns, $T_1=10\,\mu$s, $\sigma_\phi^2=0.04$: $14$ points. Tail
objectives are optimized at $\alpha\in\{0.20,0.05,0.01\}$ and every control is
evaluated at $\alpha\in\{0.5,0.2,0.05,0.01,2\times10^{-3},5\times10^{-4},
10^{-4}\}$. The fairness protocol of Sec.~\ref{sec:stage3design} is unchanged:
common warm start, budget parity, reverse warm starts, and $\Rt$ scored against
the best mean-side comparator. The comparator upgraded $13$ of the $14$ cells,
by at most $0.188$ --- a larger fraction than any earlier stage, which is
itself a statement about how flat the landscape is here.

\subsection{Pre-registered predictions and outcomes}
\begin{description}
\item[P1] The shift in the integrated excited population between the mean and
tail optima scales as $\Xi_\alpha=1/(\alpha\rho_\alpha)$.
\emph{Outcome: the estimator is mis-specified; we do not score the
prediction.} The through-origin slope of
$(n_{\mathrm{free}}-n_{\cvar})$ on $(n_{\mathrm{free}}-n_{\mathrm{mean}})$ is
$0.716$ (bootstrap CI $[0.254,3.790]$) against a predicted $3.583$, and it does
not move as the resolution gate is tightened: $0.716$ at eleven points, $0.709$
at eight, $0.660$ at five. Regression dilution would lift it; it does not.
The identity presumed that $n_{\mathrm{int}}$ is the coordinate along which the
dephasing/jump trade-off is resolved, and the spread table below shows it is
not. The sign agrees at $13$ of $14$ points.
\item[P2] $\beta$ orders $\Rt$ and $\kappa$ becomes inert.
\emph{Outcome: the second half holds, the first fails.}
$\mathrm{Spearman}(\beta,\Rt)=+0.222$ ($p=0.45$), and the same in every
subset: $+0.112$ on the twelve two-level points, $+0.168$ dropping the two
points with $\beta>10$, $+0.011$ scoring against the raw mean comparator.
$\mathrm{Spearman}(\kappa,\Rt)=+0.249$ ($p=0.39$); partial given $\beta$,
$+0.178$ ($p=0.56$). As in Sec.~\ref{sec:stage4}, the replicate scatter caps
any such correlation: run-to-run standard deviation of $\Rt$ is $0.0316$
against a total spread of $0.0402$, an attenuation ceiling of $0.62$, so the
disattenuated values are $+0.36$ and $+0.40$. Neither is significant.
\item[P3] Exceedances of the $\kappa^2$ envelope are expected, since Stage~5
breaks (A1) by construction; not scored as a failure. \emph{Outcome:}
$\max\Rt=0.030$ with six of $14$ positive. No large exceedance occurred.
\item[P4] The tail-to-mean ratio follows Eq.~(\ref{M-eq:mixture}).
\emph{Outcome: passes.} Median relative error $1.2\%$ over the twelve
two-level points, every point above $r_\alpha$, across $\beta$ from $0.098$ to
$63$. We record explicitly that the reported
$\mathrm{Spearman}(\text{pred},\text{meas})=+0.991$ is \emph{not} an
independent check --- the prediction is a monotone function of $\beta$ alone,
so it equals $\mathrm{Spearman}(\beta,\text{ratio})$ to machine precision. The
content of the test is the level agreement at no free parameters. The
four-level arm degrades to $9.9\%$ and $16.7\%$, in both cases with the
measurement below the prediction, as a third (leakage) channel would require.
\end{description}

\subsection{Weight and leverage}
The reason P2 fails is visible in how little of the jump channel the control
can reach. Over the eight controls optimized at each point --- free, mean,
three tail levels, three reverse warm starts --- the fractional spread
$(\max-\min)/\mathrm{mean}$ is a median $0.9\%$ for the jump loss $\ell_j$
(worst case $2.0\%$), a median $1.4\%$ for $n_{\mathrm{int}}$ (worst
$3.5\%$), and a median $45\%$ for the dephasing loss they are traded against
(worst $573\%$). The jump probability is pinned by the gate's own requirement
on the integrated excited population, and the cost of a relaxation event is
set by the target unitary. A tail objective multiplies the weight on that
channel by $\Xi_\alpha=3.58$; there is nothing for the extra weight to act on.

Two cells at $\sigma_\phi^2=0.16$ make the point sharply. At seed $1$ the
optimizer reaches a basin in which the linear phase response is almost
completely cancelled: $\E L_{\mathrm{deph}}$ falls by a factor of $73$ against
the seed-$0$ sibling, the mean loss improves from $2.82\times10^{-3}$ to
$1.35\times10^{-3}$, $\kappa$ rises from $0.128$ to $6.21$ because $\|A\|_F$
collapses, and $\beta$ reaches $63$. All eight controls at that point land in
the same basin, so it is a property of the problem and not of one restart.
$\Rt$ there is $-0.0008$.

\subsection{Per-point table}
$T_g$ in ns, $T_1$ in $\mu$s, ratio at $\alpha=0.05$, $\Rt$ against the best
mean comparator.
\begin{center}\footnotesize\setlength{\tabcolsep}{4pt}
\input{figures/tab_stage5.tex}
\end{center}

% =====================================================================
\section{The measurement-only predictor}
\label{sec:s13}

This section gives the construction behind Sec.~\ref{M-sec:predictor} of the
main text, the protocol correction it forced, and the gating rules used to
score it.

\subsection{Construction}

Let $\{\epsilon_k\}_{k=1}^{K}$ be the error coordinates of the test ensemble
under a control $u$, with weights $w_k$, losses
$L_k=\epsilon_k^{\mathsf T}W\epsilon_k$, and let
$T_\alpha=\{k: L_k\ge \mathrm{VaR}_\alpha\}$ carry total weight $\alpha$. The
tail-conditional displacement is
\begin{equation}
d_\alpha=\frac{1}{\alpha}\sum_{k\in T_\alpha} w_k\epsilon_k-\sum_k w_k\epsilon_k ,
\tag{S-disp}
\end{equation}
and the object of interest is the change in $\cvar_\alpha$ under a common
offset $\epsilon_k\mapsto\epsilon_k-t\,d_\alpha$ applied to every draw.

Three corrections separate a usable predictor from the naive one.

\emph{(i) Curvature.} The naive Hessian of $\cvar_\alpha$ in the offset is
$2W$, which ignores that the membership of $T_\alpha$ changes as the offset
moves. Writing $m_c$ for the tail mass and $\Sigma_\tau$ for the covariance
of the error coordinates in a narrow band around the VaR level, with
$\rho_\tau$ the density of loss there, the correct curvature is
\begin{equation}
H=\frac{m_c}{\alpha}\,2W+\frac{4\rho_\tau}{\alpha}\,W\Sigma_\tau W .
\tag{S-hess}
\end{equation}
The band is widened adaptively until it contains enough draws to estimate
$\Sigma_\tau$, and a flag is raised when it cannot. The boundary term is not
a refinement: omitting it over-predicts the gain by $315\%$ in a brute-force
comparison on the synthetic reduced model, which is how the error was found.

\emph{(ii) Reachability.} The offset must lie in
$\mathcal{B}=\operatorname{range}(\partial\epsilon/\partial c)$, computed by
finite differences on the control coefficients and whitened in the $W$
metric. On the real simulator this projector has rank $2$ of $3$ at every
point---one error coordinate is unreachable---though the measured
displacement already lies in $\mathcal{B}$ to machine precision
($\mathrm{reach\ frac}=1.000$), so the rank deficiency is not the source of
the calibration constant.

\emph{(iii) Step length.} Rather than trusting the quadratic, the step is set
by a line search of $\cvar_\alpha$ along $d_\alpha$ over a seven-decade
two-sided grid with local refinement. Because the search and the evaluation
would otherwise use the same draws, the search is cross-fitted. The fold
split is mixture-aware: splitting naively puts a rare atom entirely in one
fold, leaving that fold with no heavy tail and a spurious gain, so responsive
draws are split while inert branches are kept in both folds with renormalized
weights.

We report three nested statistics per point. $R_{\mathrm{pred}}^{\mathrm{simple}}$
uses $H=2W$ and no line search; $R_{\mathrm{pred}}^{\mathrm{quad}}$ uses
Eq.~(S-hess); and $R_{\mathrm{pred}}$ adds the cross-fitted line search.
Against a brute-force re-optimization on the synthetic model their relative
errors are $+315\%$, $\{16,77,47\}\%$ and $\{0.6,3.3,4.4\}\%$ at
$\kappa=\{0.1,0.25,0.5\}$. Only the last is used for the calibration in the
main text.

\subsection{Resolution floor}

A displacement estimated from $\alpha K$ draws has sampling scatter even when
the true gain is zero, so every point carries a floor $R_{\min}$ built from
the band statistics and the effective tail count, and a point is called
resolved when $R_{\mathrm{pred}}>R_{\min}$. On $24$ independent null draw sets
the floor matches the empirical scatter to a ratio of $0.75$. The floor
drives the test-ensemble size: the working rule is
$\alpha_{\min}\,n_{\mathrm{test}}\gtrsim400$, which with
$\alpha_{\min}=0.025$ gives $n_{\mathrm{test}}=50\,000$ and is enforced
rather than advised. The training ensemble obeys a separate rule,
$\alpha_{\min}K_{\mathrm{train}}\ge8$, because the CVaR objective's gradient
is built only from its own training tail; at the default
$K_{\mathrm{train}}=128$ that is $3.2$ draws at $\alpha=0.025$. One matched
ensemble is shared by the mean, CVaR and reverse runs at a point, since
giving the CVaR arm more data than its comparator would break the budget
parity the fairness protocol rests on.

\subsection{The shared basin pool}

The comparator of Sec.~\ref{M-sec:equiv}---best of the mean control and the
reverse warm starts---is not budget-neutral when the mean multistart has not
converged. Using the criterion that no other control at a point beats the
mean control on its own objective by more than $5\%$, only $6$ of $16$ cells
were converged, the worst at a ratio of $0.250$. Choosing the comparator by
lowest mean loss rather than lowest $\cvar$ gives numerically identical
results at every point, so this is not selection on the outcome variable; it
is an unequal restart budget.

The pooled protocol collects every control any arm found at a point
(nine, deduplicating to three or four), and runs each objective from each
start, selecting each arm's winner on the validation ensemble by that arm's
own risk measure. Table~\ref{tab:s13pool} shows the effect on the replicate
scatter. Because every start is warm, the pooled pass costs roughly
$\mathrm{len(pool)}\times5$ single optimizations against the $36$ the
original point cost.

\subsection{Gating}

Two gates are applied and both are reported alongside the ungated statistics.

For the direction test, a point enters if it is resolved, if the CVaR arm did
not find a better mean optimum than the mean control (mean-loss ratio
$\ge0.95$), and if $|R^{\mathrm{raw}}|>0.005$. The first two conditions say
that the CVaR control is a local displacement of the mean control rather than
a different basin; the third says there is a displacement at all. This
removes $4$ to $6$ of $16$ points depending on $\alpha$. The pass criterion
is significant anti-alignment---no point with $\cos<0$ whose null probability
is below $0.05$---rather than a minimum cosine, since in three dimensions
$\cos$ is uniform on $[-1,1]$ under the null and a single low value carries
almost no information.

For the level test, no gating is applied: all $16$ points enter at every
$\alpha$, and the same-basin and raw-comparator variants are reported beside
the pooled one.

\subsection{The two leverage arms}

The relaxation arm is analytic. Stage~5 never sampled jumps; it used the
exact no-jump/one-jump decomposition, recording $(p,\ell_j,\E L_c)$ and hence
$\beta$ per cell. The reconstruction appends one atom of weight $p$ carrying
loss $\ell_j$, with its error coordinates pinned at $\epsilon_0$ so that it
contributes nothing to the displacement, and rescales $\ell_j$ so that the
re-derived cell's $\beta$ matches the recorded one---$\beta$, not the
absolute jump loss, being the invariant Proposition~\ref{M-prop:mixture} is
written in. Two cells at $\kappa=7.4$ are excluded from the identity test as
far outside the perturbative regime; they satisfy it anyway. We note that
cells sharing $(T_g,\sigma_\phi^2,\text{seed},\text{levels})$ but differing
in $T_1$ produce identical re-derived controls, since $T_1$ does not enter
the dephasing bulk---so the $14$ rows are $7$ distinct bulk controls paired
with two relaxation rates, and the effective count for any statement about
the bulk is $7$.

The fluctuator arm installs a two-state telegraph process in the simulator,
so that the optimizer, the losses and the propagators all see it, and sweeps
two axes: excursion amplitude at fixed rarity, and rarity at fixed amplitude.
At each point the predictor is evaluated twice---once on the mixture and once
on the same control restricted to the non-fluctuator draws---so the
suppression ratio is taken at one control and one $\kappa$ and is not
confounded by the fact that amplitude drives $\kappa$ and rarity in opposite
directions. A rank correlation between $\Lambda$ and $R_{\mathrm{pred}}$
across an amplitude sweep is $-1$ by construction and is reported only to
make that visible.

\subsection{Validation}

Twelve checks run without the simulator and are required to pass before any
point is computed: exact atomic $\cvar$ at four $\alpha$; the vanishing of
the debiased displacement in the linear-Gaussian limit and the $1/K$ decay of
the raw bias; the closed form against brute-force re-optimization at three
$\kappa$; the predicted direction; the $\kappa^2$ onset; two leverage checks
against brute force; the floor against empirical null scatter; the two
closed-form identities for $\rho_\alpha$; and the mixture identity of
Eq.~\eqref{M-eq:suppression} against its analytic form, which agrees to $2.4\%$
and reproduces the multiplicative relation to $1.0\%$ on the synthetic model.

\begin{table}[t]
\caption{Replicate scatter in the measured gain between the two seeds, before
and after equalizing the restart budget. The pooled protocol removes almost
all of it, which is why the slope in Table~\ref{M-tab:stage6} is a calibration
rather than noise.}
\label{tab:s13pool}
\begin{ruledtabular}
\begin{tabular}{lcccc}
$\alpha$ & 0.20 & 0.10 & 0.05 & 0.025 \\ \hline
best-of-comparator & 0.272 & 0.249 & 0.178 & 0.136 \\
shared basin pool  & 0.010 & 0.016 & 0.033 & 0.019 \\
reduction factor   & 26.8  & 15.2  & 5.4   & 7.1   \\
\end{tabular}
\end{ruledtabular}
\end{table}

% =====================================================================
\section{Per-point results}
\label{sec:fulltable}
\subsection{All validated comparisons (E1--E8)}
Every cell is re-evaluated in the full Lindblad model on the test ensemble.
$\Rt$ is against the best mean control; ``same-budget'' is against the
multistart mean control alone; ``upgraded'' marks cells in which a reverse
start improved the comparator by more than 0.1\%. $T_g$ in ns,
$\sigma_\delta/2\pi$ in MHz. Point labels: \emph{Tg}, gate-time sweep;
\emph{2lv}, two levels; \emph{sig}, phase-variance sweep.
\begin{center}\scriptsize\setlength{\tabcolsep}{3pt}
\begin{longtable}{@{}llrrrrrl@{}}
\toprule
Source & Point & $T_g$ & $\sigma_\delta/2\pi$ & $\alpha$ & $R_{\rm tail}$ vs best mean [95\% CI] & same-budget & comparator\\
\midrule
\endhead
2lvl\_long & Tg-2lv 00 & 20.8 & 2.50 & 0.05 & $-0.007$ [-0.021, +0.009] & $+0.098$ & upgraded\\
2lvl\_long & Tg-2lv 00 & 20.8 & 2.50 & 0.2 & $-0.082$ [-0.094, -0.071] & $+0.026$ & upgraded\\
2lvl\_long & Tg-2lv 01 & 32.5 & 2.50 & 0.05 & $+0.009$ [-0.015, +0.033] & $+0.009$ & multistart\\
2lvl\_long & Tg-2lv 01 & 32.5 & 2.50 & 0.2 & $+0.009$ [+0.003, +0.015] & $+0.009$ & multistart\\
amp\_only & Tg 00 & 12.7 & 0.00 & 0.05 & $-0.013$ [-0.018, -0.007] & $-0.007$ & upgraded\\
amp\_only & Tg 00 & 12.7 & 0.00 & 0.2 & $-0.008$ [-0.011, -0.006] & $-0.004$ & upgraded\\
amp\_only & Tg 01 & 15.4 & 0.00 & 0.05 & $-0.166$ [-0.179, -0.154] & $-0.007$ & upgraded\\
amp\_only & Tg 01 & 15.4 & 0.00 & 0.2 & $-0.001$ [-0.004, +0.001] & $+0.112$ & upgraded\\
amp\_only & Tg 02 & 20.8 & 0.00 & 0.05 & $-0.001$ [-0.001, -0.001] & $-0.002$ & multistart\\
amp\_only & Tg 02 & 20.8 & 0.00 & 0.2 & $-0.001$ [-0.001, -0.001] & $-0.002$ & multistart\\
collapse\_T20 & sig-2lv 00 & 20.0 & 0.56 & 0.05 & $-0.001$ [-0.003, +0.001] & $+0.057$ & upgraded\\
collapse\_T20 & sig-2lv 00 & 20.0 & 0.56 & 0.2 & $-0.043$ [-0.047, -0.040] & $-0.000$ & upgraded\\
collapse\_T20 & sig-2lv 01 & 20.0 & 1.59 & 0.05 & $+0.003$ [+0.002, +0.005] & $+0.003$ & multistart\\
collapse\_T20 & sig-2lv 01 & 20.0 & 1.59 & 0.2 & $+0.002$ [-0.004, +0.009] & $+0.002$ & multistart\\
collapse\_T20 & sig-2lv 02 & 20.0 & 3.18 & 0.05 & $+0.021$ [+0.018, +0.024] & $+0.021$ & multistart\\
collapse\_T20 & sig-2lv 02 & 20.0 & 3.18 & 0.2 & $+0.010$ [+0.002, +0.019] & $+0.010$ & multistart\\
collapse\_T20 & sig-2lv 03 & 20.0 & 5.03 & 0.05 & $+0.056$ [+0.049, +0.062] & $+0.067$ & upgraded\\
collapse\_T20 & sig-2lv 03 & 20.0 & 5.03 & 0.2 & $+0.013$ [+0.003, +0.024] & $+0.026$ & upgraded\\
collapse\_T20 & sig-2lv 04 & 20.0 & 10.07 & 0.05 & $+0.292$ [+0.247, +0.335] & $+0.773$ & upgraded\\
collapse\_T20 & sig-2lv 04 & 20.0 & 10.07 & 0.2 & $-3.556$ [-4.292, -3.023] & $+0.105$ & upgraded\\
collapse\_T8 & sig-2lv 00 & 8.0 & 1.41 & 0.05 & $-0.002$ [-0.005, +0.001] & $+0.130$ & upgraded\\
collapse\_T8 & sig-2lv 00 & 8.0 & 1.41 & 0.2 & $-0.119$ [-0.126, -0.111] & $+0.000$ & upgraded\\
collapse\_T8 & sig-2lv 01 & 8.0 & 3.98 & 0.05 & $+0.001$ [-0.004, +0.005] & $+0.001$ & multistart\\
collapse\_T8 & sig-2lv 01 & 8.0 & 3.98 & 0.2 & $+0.002$ [-0.006, +0.011] & $+0.002$ & multistart\\
collapse\_T8 & sig-2lv 02 & 8.0 & 7.96 & 0.05 & $+0.018$ [+0.013, +0.022] & $+0.018$ & multistart\\
collapse\_T8 & sig-2lv 02 & 8.0 & 7.96 & 0.2 & $+0.010$ [+0.001, +0.019] & $+0.010$ & multistart\\
collapse\_T8 & sig-2lv 03 & 8.0 & 12.57 & 0.05 & $+0.051$ [+0.045, +0.056] & $+0.051$ & multistart\\
collapse\_T8 & sig-2lv 03 & 8.0 & 12.57 & 0.2 & $+0.024$ [+0.014, +0.034] & $+0.024$ & multistart\\
collapse\_T8 & sig-2lv 04 & 8.0 & 25.17 & 0.05 & $+0.187$ [+0.142, +0.224] & $+0.187$ & multistart\\
collapse\_T8 & sig-2lv 04 & 8.0 & 25.17 & 0.2 & $+0.095$ [+0.079, +0.108] & $+0.095$ & multistart\\
det\_only & Tg 00 & 12.7 & 2.50 & 0.05 & $-0.015$ [-0.035, +0.005] & $-0.011$ & upgraded\\
det\_only & Tg 00 & 12.7 & 2.50 & 0.2 & $-0.001$ [-0.010, +0.009] & $+0.003$ & upgraded\\
det\_only & Tg 01 & 15.4 & 2.50 & 0.05 & $+0.074$ [+0.058, +0.086] & $+0.154$ & upgraded\\
det\_only & Tg 01 & 15.4 & 2.50 & 0.2 & $+0.034$ [+0.019, +0.047] & $+0.111$ & upgraded\\
det\_only & Tg 02 & 20.8 & 2.50 & 0.05 & $+0.177$ [+0.153, +0.195] & $+0.167$ & upgraded\\
det\_only & Tg 02 & 20.8 & 2.50 & 0.2 & $+0.068$ [+0.050, +0.085] & $+0.077$ & upgraded\\
fast & audit & 12.7 & 2.50 & 0.01 & $-0.047$ [-0.160, +0.056] & $-0.038$ & upgraded\\
fast & audit & 12.7 & 2.50 & 0.025 & $-0.011$ [-0.025, +0.003] & $-0.000$ & upgraded\\
fast & audit & 12.7 & 2.50 & 0.05 & $-0.020$ [-0.043, -0.001] & $-0.010$ & upgraded\\
fast & audit & 12.7 & 2.50 & 0.1 & $-0.004$ [-0.014, +0.007] & $+0.006$ & upgraded\\
fast & audit & 12.7 & 2.50 & 0.2 & $-0.013$ [-0.025, -0.001] & $-0.004$ & upgraded\\
fast & sig-2lv 00 & 12.7 & 0.50 & 0.05 & $-0.001$ [-0.002, -0.001] & $+0.038$ & upgraded\\
fast & sig-2lv 00 & 12.7 & 0.50 & 0.2 & $-0.026$ [-0.028, -0.024] & $-0.000$ & upgraded\\
fast & sig-2lv 01 & 12.7 & 1.00 & 0.05 & $-0.002$ [-0.005, +0.001] & $-0.002$ & multistart\\
fast & sig-2lv 01 & 12.7 & 1.00 & 0.2 & $+0.000$ [-0.003, +0.003] & $+0.000$ & multistart\\
fast & sig-2lv 02 & 12.7 & 2.50 & 0.05 & $+0.002$ [-0.003, +0.007] & $+0.002$ & multistart\\
fast & sig-2lv 02 & 12.7 & 2.50 & 0.2 & $+0.003$ [-0.004, +0.009] & $+0.002$ & multistart\\
fast & sig-2lv 03 & 12.7 & 5.00 & 0.05 & $+0.029$ [+0.026, +0.032] & $+0.029$ & multistart\\
fast & sig-2lv 03 & 12.7 & 5.00 & 0.2 & $+0.013$ [+0.004, +0.022] & $+0.013$ & multistart\\
fast & sig-2lv 04 & 12.7 & 8.00 & 0.05 & $+0.077$ [+0.068, +0.085] & $+0.077$ & multistart\\
fast & sig-2lv 04 & 12.7 & 8.00 & 0.2 & $+0.034$ [+0.024, +0.045] & $+0.034$ & multistart\\
fast & sig-2lv 05 & 12.7 & 16.00 & 0.05 & $+0.275$ [+0.247, +0.298] & $+0.275$ & multistart\\
fast & sig-2lv 05 & 12.7 & 16.00 & 0.2 & $+0.135$ [+0.116, +0.152] & $+0.135$ & multistart\\
fast & Tg 00 & 6.3 & 2.50 & 0.05 & $-0.003$ [-0.007, +0.002] & $-0.002$ & upgraded\\
fast & Tg 00 & 6.3 & 2.50 & 0.2 & $-0.002$ [-0.014, +0.010] & $+0.001$ & upgraded\\
fast & Tg 01 & 7.4 & 2.50 & 0.05 & $-0.011$ [-0.024, +0.003] & $-0.011$ & multistart\\
fast & Tg 01 & 7.4 & 2.50 & 0.2 & $+0.001$ [-0.006, +0.007] & $+0.001$ & multistart\\
fast & Tg 02 & 8.5 & 2.50 & 0.05 & $-0.016$ [-0.030, -0.001] & $-0.013$ & upgraded\\
fast & Tg 02 & 8.5 & 2.50 & 0.2 & $-0.002$ [-0.011, +0.006] & $-0.000$ & upgraded\\
fast & Tg 03 & 9.9 & 2.50 & 0.05 & $-0.017$ [-0.031, -0.001] & $-0.016$ & multistart\\
fast & Tg 03 & 9.9 & 2.50 & 0.2 & $-0.002$ [-0.011, +0.008] & $-0.001$ & multistart\\
fast & Tg 04 & 11.5 & 2.50 & 0.05 & $-0.021$ [-0.041, -0.002] & $-0.021$ & multistart\\
fast & Tg 04 & 11.5 & 2.50 & 0.2 & $-0.004$ [-0.016, +0.008] & $-0.004$ & multistart\\
fast & Tg 05 & 12.7 & 2.50 & 0.05 & $-0.022$ [-0.046, -0.001] & $-0.013$ & upgraded\\
fast & Tg 05 & 12.7 & 2.50 & 0.2 & $-0.013$ [-0.026, -0.000] & $-0.006$ & upgraded\\
fast & Tg 06 & 15.4 & 2.50 & 0.05 & $+0.104$ [+0.080, +0.125] & $+0.107$ & upgraded\\
fast & Tg 06 & 15.4 & 2.50 & 0.2 & $+0.038$ [+0.027, +0.049] & $+0.042$ & upgraded\\
fast & Tg 07 & 20.8 & 2.50 & 0.05 & $+0.102$ [+0.074, +0.123] & $+0.102$ & multistart\\
fast & Tg 07 & 20.8 & 2.50 & 0.2 & $+0.069$ [+0.057, +0.079] & $+0.069$ & multistart\\
fast & Tg 08 & 32.5 & 2.50 & 0.05 & $+0.064$ [+0.035, +0.090] & $+0.079$ & upgraded\\
fast & Tg 08 & 32.5 & 2.50 & 0.2 & $-0.022$ [-0.034, -0.009] & $+0.024$ & upgraded\\
fast & Tg-2lv 00 & 6.3 & 2.50 & 0.05 & $-0.026$ [-0.053, -0.000] & $-0.026$ & multistart\\
fast & Tg-2lv 00 & 6.3 & 2.50 & 0.2 & $-0.000$ [-0.003, +0.003] & $-0.000$ & multistart\\
fast & Tg-2lv 01 & 12.7 & 2.50 & 0.05 & $-0.026$ [-0.053, +0.003] & $-0.026$ & multistart\\
fast & Tg-2lv 01 & 12.7 & 2.50 & 0.2 & $-0.005$ [-0.011, +0.001] & $-0.005$ & multistart\\
fast & Tg-2lv 02 & 15.4 & 2.50 & 0.05 & $-0.012$ [-0.030, +0.004] & $-0.012$ & multistart\\
fast & Tg-2lv 02 & 15.4 & 2.50 & 0.2 & $-0.000$ [-0.003, +0.004] & $-0.000$ & multistart\\
seed1 & audit & 12.7 & 2.50 & 0.01 & $-0.145$ [-0.176, -0.080] & $-0.137$ & upgraded\\
seed1 & audit & 12.7 & 2.50 & 0.025 & $-0.047$ [-0.082, -0.010] & $-0.041$ & upgraded\\
seed1 & audit & 12.7 & 2.50 & 0.05 & $-0.041$ [-0.075, -0.011] & $-0.034$ & upgraded\\
seed1 & audit & 12.7 & 2.50 & 0.1 & $-0.044$ [-0.072, -0.017] & $-0.038$ & upgraded\\
seed1 & audit & 12.7 & 2.50 & 0.2 & $-0.015$ [-0.020, -0.011] & $-0.010$ & upgraded\\
seed1 & sig-2lv 00 & 12.7 & 5.00 & 0.05 & $-0.396$ [-0.430, -0.362] & $+0.050$ & upgraded\\
seed1 & sig-2lv 00 & 12.7 & 5.00 & 0.2 & $+0.005$ [+0.000, +0.010] & $+0.318$ & upgraded\\
seed1 & sig-2lv 01 & 12.7 & 16.00 & 0.05 & $+0.271$ [+0.219, +0.316] & $+0.271$ & multistart\\
seed1 & sig-2lv 01 & 12.7 & 16.00 & 0.2 & $+0.097$ [+0.059, +0.127] & $+0.097$ & multistart\\
seed2 & audit & 12.7 & 2.50 & 0.01 & $+0.014$ [+0.001, +0.029] & $+0.016$ & upgraded\\
seed2 & audit & 12.7 & 2.50 & 0.025 & $+0.010$ [-0.017, +0.033] & $+0.013$ & upgraded\\
seed2 & audit & 12.7 & 2.50 & 0.05 & $+0.015$ [+0.000, +0.029] & $+0.018$ & upgraded\\
seed2 & audit & 12.7 & 2.50 & 0.1 & $+0.003$ [-0.008, +0.012] & $+0.005$ & upgraded\\
seed2 & audit & 12.7 & 2.50 & 0.2 & $-0.004$ [-0.012, +0.003] & $-0.001$ & upgraded\\
seed2 & sig-2lv 00 & 12.7 & 5.00 & 0.05 & $+0.036$ [+0.024, +0.046] & $+0.036$ & multistart\\
seed2 & sig-2lv 00 & 12.7 & 5.00 & 0.2 & $+0.011$ [+0.008, +0.014] & $+0.011$ & multistart\\
seed2 & sig-2lv 01 & 12.7 & 16.00 & 0.05 & $+0.314$ [+0.279, +0.341] & $+0.314$ & multistart\\
seed2 & sig-2lv 01 & 12.7 & 16.00 & 0.2 & $+0.142$ [+0.127, +0.154] & $+0.142$ & multistart\\
strong & Tg 00 & 15.4 & 2.50 & 0.05 & $+0.100$ [+0.077, +0.122] & $+0.100$ & multistart\\
strong & Tg 00 & 15.4 & 2.50 & 0.2 & $+0.030$ [+0.020, +0.039] & $+0.030$ & multistart\\
strong & Tg 01 & 20.8 & 2.50 & 0.05 & $+0.117$ [+0.092, +0.138] & $+0.124$ & upgraded\\
strong & Tg 01 & 20.8 & 2.50 & 0.2 & $+0.060$ [+0.046, +0.072] & $+0.066$ & upgraded\\
strong & Tg 02 & 32.5 & 2.50 & 0.05 & $+0.059$ [+0.036, +0.080] & $+0.059$ & multistart\\
strong & Tg 02 & 32.5 & 2.50 & 0.2 & $+0.027$ [+0.017, +0.037] & $+0.027$ & multistart\\
strong\_2lvl & Tg-2lv 00 & 15.4 & 2.50 & 0.05 & $-0.008$ [-0.038, +0.022] & $+0.076$ & upgraded\\
strong\_2lvl & Tg-2lv 00 & 15.4 & 2.50 & 0.2 & $-0.052$ [-0.069, -0.033] & $+0.000$ & upgraded\\
strong\_2lvl & Tg-2lv 01 & 20.8 & 2.50 & 0.05 & $-0.021$ [-0.043, -0.001] & $-0.021$ & multistart\\
strong\_2lvl & Tg-2lv 01 & 20.8 & 2.50 & 0.2 & $+0.000$ [-0.001, +0.002] & $+0.000$ & multistart\\
strong\_2lvl & Tg-2lv 02 & 32.5 & 2.50 & 0.05 & $-0.004$ [-0.010, +0.001] & $-0.000$ & upgraded\\
strong\_2lvl & Tg-2lv 02 & 32.5 & 2.50 & 0.2 & $-0.000$ [-0.007, +0.006] & $+0.004$ & upgraded\\
\bottomrule
\end{longtable}
\end{center}

\subsection{Dominance test (E9)}
Closed-system statistics over each candidate pool (fresh 4000-draw ensemble)
and the full-Lindblad $\Rt(0.05)$ of the standard comparison against the best
mean control. $\kappa$: second-order ratio of the mean control. $\tau$:
Kendall $\tau$ between mean and $\cvar_{0.05}$ rankings. Rev.: fraction of
confirmed rank reversals. Dom.: fraction of mean-separated pairs crossing for
$q\ge0.5$. Ratio: median $\cvar_{0.05}/$mean (prediction $5.58$).
$R_{\rm pool}$: cross-fitted pool regret. $^{\circ}$Comparator upgraded by a
reverse start.
{\scriptsize\setlength{\tabcolsep}{3pt}
\begin{longtable}{@{}rrrrrrrrrrr@{}}
\toprule
Lev. & $T_g$ & seed & $\sigma_\phi^2$ & $\kappa$ & $\tau$ & Rev. & Dom. & Ratio & $R_{\rm pool}$ & $\Rt$ Lindblad\\
\midrule
\endhead
2 & 8 & 0 & 0.0016 & 0.01 & 0.98 & 0.000 & 0.000 & 5.60 & $+0.000$ & $-0.001$$^{\circ}$ \\
2 & 8 & 0 & 0.0064 & 0.03 & 0.95 & 0.000 & 0.000 & 5.59 & $+0.000$ & $-0.014$$^{\circ}$ \\
2 & 8 & 0 & 0.04 & 0.07 & 0.80 & 0.001 & 0.000 & 5.52 & $+0.003$ & $+0.001$$^{\circ}$ \\
2 & 8 & 0 & 0.16 & 0.13 & 0.62 & 0.095 & 0.054 & 5.37 & $+0.022$ & $+0.017$$^{\circ}$ \\
2 & 8 & 0 & 0.41 & 0.21 & 0.47 & 0.226 & 0.256 & 4.95 & $+0.053$ & $+0.052$$^{\circ}$ \\
2 & 8 & 0 & 1.6 & 0.43 & 0.62 & 0.165 & 0.269 & 9.81 & $+0.178$ & $+0.191$$^{\circ}$ \\
2 & 12.7 & 0 & 0.0016 & 0.01 & 0.94 & 0.000 & 0.000 & 5.60 & $+0.000$ & $-0.001$$^{\circ}$ \\
2 & 12.7 & 1 & 0.0016 & 0.01 & 0.90 & 0.000 & 0.000 & 5.59 & $+0.000$ & $-0.040$$^{\circ}$ \\
2 & 12.7 & 2 & 0.0016 & 0.01 & 0.97 & 0.000 & 0.000 & 5.49 & $+0.000$ & $-0.032$$^{\circ}$ \\
2 & 12.7 & 0 & 0.0064 & 0.03 & 0.95 & 0.000 & 0.000 & 5.59 & $+0.000$ & $-0.002$ \\
2 & 12.7 & 1 & 0.0064 & 0.03 & 0.95 & 0.000 & 0.000 & 5.59 & $+0.001$ & $-0.013$ \\
2 & 12.7 & 2 & 0.0064 & 0.03 & 0.95 & 0.000 & 0.000 & 5.49 & $+0.001$ & $+0.001$ \\
2 & 12.7 & 0 & 0.04 & 0.08 & 0.83 & 0.001 & 0.000 & 5.51 & $+0.002$ & $+0.003$$^{\circ}$ \\
2 & 12.7 & 1 & 0.04 & 0.08 & 0.87 & 0.008 & 0.000 & 5.52 & $+0.003$ & $-0.040$$^{\circ}$ \\
2 & 12.7 & 2 & 0.04 & 0.08 & 0.86 & 0.005 & 0.000 & 5.46 & $-0.001$ & $-0.001$$^{\circ}$ \\
2 & 12.7 & 0 & 0.16 & 0.15 & 0.65 & 0.125 & 0.040 & 5.50 & $+0.032$ & $+0.027$$^{\circ}$ \\
2 & 12.7 & 1 & 0.16 & 0.15 & 0.68 & 0.123 & 0.050 & 5.37 & $+0.028$ & $+0.001$$^{\circ}$ \\
2 & 12.7 & 2 & 0.16 & 0.15 & 0.63 & 0.105 & 0.052 & 5.38 & $+0.034$ & $+0.033$ \\
2 & 12.7 & 0 & 0.41 & 0.25 & 0.59 & 0.190 & 0.192 & 5.36 & $+0.079$ & $+0.072$ \\
2 & 12.7 & 1 & 0.41 & 0.25 & 0.55 & 0.238 & 0.273 & 5.46 & $+0.072$ & $+0.044$$^{\circ}$ \\
2 & 12.7 & 2 & 0.41 & 0.25 & 0.57 & 0.198 & 0.156 & 5.29 & $+0.090$ & $+0.089$$^{\circ}$ \\
2 & 12.7 & 0 & 1.6 & 0.52 & 0.29 & 0.364 & 0.468 & 4.82 & $+0.309$ & $+0.274$$^{\circ}$ \\
2 & 12.7 & 1 & 1.6 & 0.50 & 0.34 & 0.315 & 0.361 & 4.85 & $+0.292$ & $+0.276$$^{\circ}$ \\
2 & 12.7 & 2 & 1.6 & 0.52 & 0.10 & 0.452 & 0.607 & 4.61 & $+0.309$ & $+0.290$ \\
2 & 20 & 0 & 0.0016 & 0.01 & 0.98 & 0.000 & 0.000 & 5.60 & $+0.000$ & $-0.011$$^{\circ}$ \\
2 & 20 & 0 & 0.0064 & 0.03 & 0.95 & 0.000 & 0.000 & 5.59 & $-0.002$ & $-0.001$ \\
2 & 20 & 0 & 0.04 & 0.07 & 0.83 & 0.002 & 0.000 & 5.58 & $+0.003$ & $+0.001$$^{\circ}$ \\
2 & 20 & 0 & 0.16 & 0.14 & 0.56 & 0.135 & 0.027 & 5.54 & $+0.028$ & $+0.022$$^{\circ}$ \\
2 & 20 & 0 & 0.41 & 0.23 & 0.58 & 0.167 & 0.177 & 5.43 & $+0.071$ & $+0.062$$^{\circ}$ \\
2 & 20 & 0 & 1.6 & 0.48 & 0.10 & 0.473 & 0.675 & 11.04 & $+0.342$ & $+0.353$$^{\circ}$ \\
4 & 12.7 & 0 & 0.0016 & 0.01 & 0.92 & 0.009 & 0.013 & 1.88 & $+0.001$ & $-0.003$$^{\circ}$ \\
4 & 12.7 & 0 & 0.0064 & 0.02 & 0.90 & 0.010 & 0.015 & 3.06 & $+0.002$ & $-0.013$$^{\circ}$ \\
4 & 12.7 & 0 & 0.04 & 0.04 & 0.85 & 0.018 & 0.007 & 4.73 & $+0.014$ & $-0.008$ \\
4 & 12.7 & 0 & 0.16 & 0.08 & 0.71 & 0.068 & 0.055 & 5.17 & $+0.040$ & $+0.027$ \\
4 & 12.7 & 0 & 0.41 & 0.12 & 0.57 & 0.175 & 0.229 & 5.04 & $+0.081$ & $+0.068$ \\
4 & 12.7 & 0 & 1.6 & 0.26 & 0.55 & 0.165 & 0.232 & 4.15 & $+0.048$ & $+0.042$$^{\circ}$ \\
4 & 20.8 & 0 & 0.0016 & 0.01 & 0.96 & 0.002 & 0.000 & 1.75 & $+0.000$ & $-0.002$$^{\circ}$ \\
4 & 20.8 & 0 & 0.0064 & 0.02 & 0.95 & 0.001 & 0.000 & 2.59 & $+0.001$ & $-0.000$$^{\circ}$ \\
4 & 20.8 & 0 & 0.04 & 0.06 & 0.74 & 0.046 & 0.026 & 3.32 & $+0.029$ & $+0.016$$^{\circ}$ \\
4 & 20.8 & 0 & 0.16 & 0.12 & 0.49 & 0.220 & 0.266 & 4.67 & $+0.123$ & $+0.104$$^{\circ}$ \\
4 & 20.8 & 0 & 0.41 & 0.20 & 0.20 & 0.380 & 0.527 & 5.16 & $+0.278$ & $+0.274$ \\
4 & 20.8 & 0 & 1.6 & 0.38 & 0.19 & 0.397 & 0.579 & 4.93 & $+0.231$ & $+0.240$ \\
\bottomrule
\end{longtable}}
% =====================================================================
% =====================================================================
\section{Reproducibility}
\label{sec:repro}

\subsection{Code}

All code is plain Python 3 with NumPy and SciPy (plus Matplotlib for the
figures). Every step writes machine-readable output (JSON, CSV and NPZ with
the full control vectors and per-draw test losses), so every table and figure
can be regenerated from saved data without re-running an optimization.

\begin{center}\small
\begin{tabular}{@{}p{0.27\textwidth}p{0.66\textwidth}@{}}
\toprule
File & Role\\
\midrule
\texttt{fade\_realistic.py} & transmon model, thermal state, Butterworth filter matrix, Hann seed pulse, worker pool\\
\texttt{stage0.py} & per-draw full-Lindblad worker with exact (Van Loan) gradients\\
\texttt{stage0\_v2.py} & configuration, noise ensembles (with correlation), CVaR weights, bounded pulse basis, optimizer, fairness protocol, paired bootstrap, Stage-0 steps\\
\texttt{fast\_viper.py} & closed-system Daleckii--Krein simulator, continuation starts, self-test against Lindblad, full-Lindblad \texttt{validate} step with the best-mean comparator\\
\texttt{stage1.py} & Stage-1 \texttt{channels} runs and the \texttt{diagnose} step (error rotation vectors, loss-quadratic fits, error clouds)\\
\texttt{run\_final.sh} & the final Stage-0 batch (collapse test, replication, strong search, validation)\\
\texttt{make\_figures.py} & Stage-0 figures and the Sec.~\ref{sec:fulltable} table\\
\texttt{make\_stage1\_figures.py} & $\kappa$, reduced model, offsets, independence test\\
\texttt{make\_paper\_extras.py} & explanatory figures, pulses, loss distributions, error clouds, simulator check\\
\texttt{stage2.py}, \texttt{run\_stage2.sh} & Stage-2 candidate pools and dominance statistics\\
\texttt{make\_stage2\_figures.py} & Stage-2 figures, revised dominance test, $\kappa$, table\\
\texttt{make\_kappa\_bound.py} & the $\kappa^2$ law: closed-form coefficients, the scalar and random-geometry checks, and the comparison with every diagnosed control\\
\texttt{stage3.py}, \texttt{run\_stage3.sh} & the controlled $\kappa$ experiment, leakage-aware diagnosis and the generality grid\\
\texttt{stage3\_recomp.py} & applies the best-mean comparator to the saved closed-system losses\\
\texttt{make\_stage3\_figures.py} & Stage-3 figures and the partial-correlation analysis\\
\bottomrule
\end{tabular}
\end{center}

\paragraph*{Built-in checks.} \texttt{stage0\_v2.py selftest} and
\texttt{fast\_viper.py selftest} must pass before any production run. They
check gradients against finite differences, the closed-system model against
the Lindblad code with dissipation switched off ($10^{-7}$), that dissipation
acts as a near-constant shift, and that the CVaR weights reproduce the exact
estimator. \texttt{stage1.py} refuses to run with an older version of the
ensemble code that lacks the correlation parameter.

\subsection{Commands}

Stage~0 (from the project folder, 8 worker processes):
\begin{verbatim}
export FADE_WORKERS=8
python fast_viper.py selftest                 # default outdir results_fast
python fast_viper.py characterize             # leakage onset
python fast_viper.py audit --Tg-audit 12.7
python fast_viper.py sweep --Tg 6.337 7.351 8.529 9.894 11.479 12.7 \
       15.449 20.793 32.467                   # E2, 4 levels
python fast_viper.py sweep --levels 2 --Tg 6.337 12.7 15.449
python fast_viper.py sweep --levels 2 --Tg 20.793 32.467 \
       --outdir results_2lvl_long
python fast_viper.py sweep --Tg 12.7 15.449 20.793 --sigma-d-mhz 0 \
       --outdir results_amp_only              # E3
python fast_viper.py sweep --Tg 12.7 15.449 20.793 --sigma-a 0 \
       --outdir results_det_only
python fast_viper.py sigma --levels 2 --Tg-sigma 12.7 \
       --sigma-grid 0.5 1 2.5 5 8 16          # E4 at 12.7 ns
python fast_viper.py sweep --levels 2 --Tg 15.449 20.793 32.467 --restarts 20 \
       --maxiter 1000 --outdir results_strong_2lvl
bash run_final.sh 2>&1 | tee final_log.txt   # collapse, seeds, strong search
\end{verbatim}
(\texttt{fast\_viper.py} passes every step except \texttt{selftest} and
\texttt{validate} to \texttt{stage0\_v2.py} with the fast simulator patched
in.) \texttt{run\_final.sh}
contains, among others,
\begin{verbatim}
python fast_viper.py sigma --levels 2 --Tg-sigma 8 \
       --sigma-grid 1.41 3.98 7.96 12.57 25.17 --outdir results_collapse_T8
python fast_viper.py sigma --levels 2 --Tg-sigma 20 \
       --sigma-grid 0.56 1.59 3.18 5.03 10.07 --outdir results_collapse_T20
python fast_viper.py audit --Tg-audit 12.7 --seed $s --outdir results_seed$s
python fast_viper.py sweep --Tg 15.449 20.793 32.467 --restarts 20 \
       --maxiter 1000 --outdir results_strong
python fast_viper.py validate --outdir <folder>     # for every folder
\end{verbatim}

Stage~1 (2 levels, $T_g=12.7$~ns unless stated):
\begin{verbatim}
python stage1.py channels --sigma-a-grid 0.01 0.02 0.04 0.08 0.127
python stage1.py channels --sigma-a-grid 0.127 \
       --rho-grid -0.99 -0.9 -0.5 0.5 0.9 0.99
python stage1.py channels --sigma-d-mhz 0 --sigma-a-grid 0.127
python stage1.py channels --sigma-a-grid 0.08 --seed 1     # and --seed 2
python stage1.py channels --sigma-a-grid 0.127 --rho-grid 0.99 --seed 1
python fast_viper.py validate --outdir results_stage1
python stage1.py diagnose results_stage1 results_fast ...
\end{verbatim}
\sloppy
The prospective test (E8) used the same \texttt{channels} step at
$\rho=0.99$ with (noise $\times0.25$)
\texttt{--sigma-d-mhz 0.625
--sigma-a-grid 0.03175}, (noise $\times0.5$) \texttt{--sigma-d-mhz 1.25
--sigma-a-grid 0.0635}, and at matched phase variance
\texttt{--Tg 8 --sigma-d-mhz 3.97 --sigma-a-grid 0.127} and
\texttt{--Tg 20 --sigma-d-mhz 1.5875 --sigma-a-grid 0.127}, each followed by
\texttt{validate} and \texttt{diagnose}. The restart seeds of a
\texttt{channels} point depend on its position in the grid of that call and on
\texttt{--seed}; they are recorded in each point's JSON output.

Stage~2 (dominance test, overnight):
\begin{verbatim}
python stage2.py selftest                 # synthetic checks of every statistic
bash run_stage2.sh 2>&1 | tee stage2_log.txt
python stage2.py summary                  # per-point table
\end{verbatim}
\texttt{run\_stage2.sh} runs \texttt{stage2.py run} for the three grids of
E9, then \texttt{validate}, \texttt{diagnose} and \texttt{summary}. The pool
ensemble uses seed $44+1000s$.

Figures and tables:
\begin{verbatim}
python make_figures.py              # Stage 0 figures, per-point table
python make_stage1_figures.py       # kappa, model, offsets, independence
python make_paper_extras.py         # explanatory and Stage-1 detail figures
python make_stage2_figures.py       # Stage-2 figures, table, statistics
\end{verbatim}

\subsection{Seeds}

Training, validation and test ensembles use seeds 11, 22 and 33, shifted by
$1000s$ for replication seed $s$. Restart perturbations use a separate seed
per objective and point. Bootstraps use 2000 resamples with seed 7 in the
full-Lindblad validation. The reduced-model Monte Carlo uses $6\times10^4$
draws with seed 1, and the pre-compensation toy $2\times10^6$ draws with
seed 0.

% =====================================================================
\begin{thebibliography}{99}
\bibitem{Rockafellar2000} R.~T.~Rockafellar and S.~Uryasev, Optimization of conditional value-at-risk, J.\ Risk \textbf{2}(3), 21 (2000).
\bibitem{Acerbi2002} C.~Acerbi and D.~Tasche, On the coherence of expected shortfall, J.\ Bank.\ Finance \textbf{26}, 1487 (2002).
\bibitem{Nielsen2002} M.~A.~Nielsen, A simple formula for the average gate fidelity of a quantum dynamical operation, Phys.\ Lett.\ A \textbf{303}, 249 (2002).
\bibitem{Shaked2007} M.~Shaked and J.~G.~Shanthikumar, \emph{Stochastic Orders} (Springer, New York, 2007).
\bibitem{Marshall2011} A.~W.~Marshall, I.~Olkin, and B.~C.~Arnold, \emph{Inequalities: Theory of Majorization and Its Applications}, 2nd ed.\ (Springer, New York, 2011).
\bibitem{Koch2007} J.~Koch, T.~M.~Yu, J.~Gambetta, A.~A.~Houck, D.~I.~Schuster, J.~Majer, A.~Blais, M.~H.~Devoret, S.~M.~Girvin, and R.~J.~Schoelkopf, Charge-insensitive qubit design derived from the Cooper pair box, Phys.\ Rev.\ A \textbf{76}, 042319 (2007).
\bibitem{Krantz2019} P.~Krantz, M.~Kjaergaard, F.~Yan, T.~P.~Orlando, S.~Gustavsson, and W.~D.~Oliver, A quantum engineer's guide to superconducting qubits, Appl.\ Phys.\ Rev.\ \textbf{6}, 021318 (2019).
\bibitem{Motzoi2009} F.~Motzoi, J.~M.~Gambetta, P.~Rebentrost, and F.~K.~Wilhelm, Simple pulses for elimination of leakage in weakly nonlinear qubits, Phys.\ Rev.\ Lett.\ \textbf{103}, 110501 (2009).
\bibitem{VanLoan1978} C.~F.~Van~Loan, Computing integrals involving the matrix exponential, IEEE Trans.\ Autom.\ Control \textbf{23}, 395 (1978).
\bibitem{Higham2008} N.~J.~Higham, \emph{Functions of Matrices: Theory and Computation} (SIAM, Philadelphia, 2008).
\bibitem{Byrd1995} R.~H.~Byrd, P.~Lu, J.~Nocedal, and C.~Zhu, A limited memory algorithm for bound constrained optimization, SIAM J.\ Sci.\ Comput.\ \textbf{16}, 1190 (1995).
\bibitem{Efron1993} B.~Efron and R.~J.~Tibshirani, \emph{An Introduction to the Bootstrap} (Chapman \& Hall, New York, 1993).
\end{thebibliography}

\end{document}
